\nonstopmode{}
\documentclass[a4paper]{book}
\usepackage[times,inconsolata,hyper]{Rd}
\usepackage{makeidx}
\makeatletter\@ifl@t@r\fmtversion{2018/04/01}{}{\usepackage[utf8]{inputenc}}\makeatother
% \usepackage{graphicx} % @USE GRAPHICX@
\makeindex{}
\begin{document}
\chapter*{}
\begin{center}
{\textbf{\huge Package `radiant.model'}}
\par\bigskip{\large \today}
\end{center}
\ifthenelse{\boolean{Rd@use@hyper}}{\hypersetup{pdftitle = {radiant.model: Model Menu for Radiant: Business Analytics using R and Shiny}}}{}
\ifthenelse{\boolean{Rd@use@hyper}}{\hypersetup{pdfauthor = {Vincent Nijs}}}{}
\begin{description}
\raggedright{}
\item[Type]\AsIs{Package}
\item[Title]\AsIs{Model Menu for Radiant: Business Analytics using R and Shiny}
\item[Version]\AsIs{1.6.9}
\item[Date]\AsIs{2026-1-8}
\item[Description]\AsIs{The Radiant Model menu includes interfaces for linear and logistic
regression, naive Bayes, neural networks, classification and regression trees,
model evaluation, collaborative filtering, decision analysis, and simulation.
The application extends the functionality in 'radiant.data'.}
\item[Depends]\AsIs{R (>= 4.3.0), radiant.data (>= 1.6.6)}
\item[Imports]\AsIs{radiant.basics (>= 1.6.6), shiny (>= 1.8.1), nnet (>= 7.3.12),
NeuralNetTools (>= 1.5.1), sandwich (>= 2.3.4), car (>= 2.1.3),
ggplot2 (>= 3.4.2), scales (>= 1.2.1), data.tree (>= 0.7.4),
stringr (>= 1.1.0), lubridate (>= 1.7.2), tidyr (>= 0.8.2),
dplyr (>= 1.1.2), tidyselect (>= 1.2.0), rlang (>= 0.4.10),
magrittr (>= 1.5), DiagrammeR (>= 1.0.9), import (>= 1.1.0),
psych (>= 1.8.4), e1071 (>= 1.6.8), rpart (>= 4.1.11), ggrepel
(>= 0.8), broom (>= 0.7.0), patchwork (>= 1.0.0), ranger (>=
0.11.2), xgboost (>= 1.6.0.1), stringi, yaml}
\item[Suggests]\AsIs{testthat (>= 2.0.0), pkgdown (>= 1.1.0)}
\item[URL]\AsIs{}\url{https://github.com/radiant-rstats/radiant.model/}\AsIs{,
}\url{https://radiant-rstats.github.io/radiant.model/}\AsIs{,
}\url{https://radiant-rstats.github.io/docs/}\AsIs{}
\item[BugReports]\AsIs{}\url{https://github.com/radiant-rstats/radiant.model/issues/}\AsIs{}
\item[License]\AsIs{AGPL-3 | file LICENSE}
\item[LazyData]\AsIs{true}
\item[Encoding]\AsIs{UTF-8}
\item[Language]\AsIs{en-US}
\item[RoxygenNote]\AsIs{7.3.3}
\item[NeedsCompilation]\AsIs{no}
\item[Author]\AsIs{Vincent Nijs [aut, cre]}
\item[Maintainer]\AsIs{Vincent Nijs }\email{radiant@rady.ucsd.edu}\AsIs{}
\end{description}
\Rdcontents{Contents}
\HeaderA{.as\_int}{Convenience function used in "simulater"}{.as.Rul.int}
%
\begin{Description}
Convenience function used in "simulater"
\end{Description}
%
\begin{Usage}
\begin{verbatim}
.as_int(x, dataset = list())
\end{verbatim}
\end{Usage}
%
\begin{Arguments}
\begin{ldescription}
\item[\code{x}] Character vector to be converted to integer

\item[\code{dataset}] Data list
\end{ldescription}
\end{Arguments}
%
\begin{Value}
An integer vector
\end{Value}
\HeaderA{.as\_num}{Convenience function used in "simulater"}{.as.Rul.num}
%
\begin{Description}
Convenience function used in "simulater"
\end{Description}
%
\begin{Usage}
\begin{verbatim}
.as_num(x, dataset = list())
\end{verbatim}
\end{Usage}
%
\begin{Arguments}
\begin{ldescription}
\item[\code{x}] Character vector to be converted to an numeric value

\item[\code{dataset}] Data list
\end{ldescription}
\end{Arguments}
%
\begin{Value}
An numeric vector
\end{Value}
\HeaderA{auc}{Area Under the RO Curve (AUC)}{auc}
%
\begin{Description}
Area Under the RO Curve (AUC)
\end{Description}
%
\begin{Usage}
\begin{verbatim}
auc(pred, rvar, lev)
\end{verbatim}
\end{Usage}
%
\begin{Arguments}
\begin{ldescription}
\item[\code{pred}] Prediction or predictor

\item[\code{rvar}] Response variable

\item[\code{lev}] The level in the response variable defined as success
\end{ldescription}
\end{Arguments}
%
\begin{Details}
See \url{https://radiant-rstats.github.io/docs/model/evalbin.html} for an example in Radiant
\end{Details}
%
\begin{Value}
AUC statistic
\end{Value}
%
\begin{SeeAlso}
\code{\LinkA{evalbin}{evalbin}} to calculate results

\code{\LinkA{summary.evalbin}{summary.evalbin}} to summarize results

\code{\LinkA{plot.evalbin}{plot.evalbin}} to plot results
\end{SeeAlso}
%
\begin{Examples}
\begin{ExampleCode}
auc(runif(20000), dvd$buy, "yes")
auc(ifelse(dvd$buy == "yes", 1, 0), dvd$buy, "yes")
\end{ExampleCode}
\end{Examples}
\HeaderA{catalog}{Catalog sales for men's and women's apparel}{catalog}
\keyword{datasets}{catalog}
%
\begin{Description}
Catalog sales for men's and women's apparel
\end{Description}
%
\begin{Usage}
\begin{verbatim}
data(catalog)
\end{verbatim}
\end{Usage}
%
\begin{Format}
A data frame with 200 rows and 5 variables
\end{Format}
%
\begin{Details}
Description provided in attr(catalog, "description")
\end{Details}
\HeaderA{confint\_robust}{Confidence interval for robust estimators}{confint.Rul.robust}
%
\begin{Description}
Confidence interval for robust estimators
\end{Description}
%
\begin{Usage}
\begin{verbatim}
confint_robust(object, level = 0.95, dist = "norm", vcov = NULL, ...)
\end{verbatim}
\end{Usage}
%
\begin{Arguments}
\begin{ldescription}
\item[\code{object}] A fitted model object

\item[\code{level}] The confidence level required

\item[\code{dist}] Distribution to use ("norm" or "t")

\item[\code{vcov}] Covariance matrix generated by, e.g., sandwich::vcovHC

\item[\code{...}] Additional argument(s) for methods
\end{ldescription}
\end{Arguments}
%
\begin{Details}
Wrapper for confint with robust standard errors. See \url{https://stackoverflow.com/questions/3817182/vcovhc-and-confidence-interval/3820125\#3820125}
\end{Details}
\HeaderA{confusion}{Confusion matrix}{confusion}
%
\begin{Description}
Confusion matrix
\end{Description}
%
\begin{Usage}
\begin{verbatim}
confusion(
  dataset,
  pred,
  rvar,
  lev = "",
  cost = 1,
  margin = 2,
  scale = 1,
  train = "All",
  data_filter = "",
  arr = "",
  rows = NULL,
  envir = parent.frame(),
  ...
)
\end{verbatim}
\end{Usage}
%
\begin{Arguments}
\begin{ldescription}
\item[\code{dataset}] Dataset

\item[\code{pred}] Predictions or predictors

\item[\code{rvar}] Response variable

\item[\code{lev}] The level in the response variable defined as success

\item[\code{cost}] Cost for each connection (e.g., email or mailing)

\item[\code{margin}] Margin on each customer purchase

\item[\code{scale}] Scaling factor to apply to calculations

\item[\code{train}] Use data from training ("Training"), test ("Test"), both ("Both"), or all data ("All") to evaluate model evalbin

\item[\code{data\_filter}] Expression entered in, e.g., Data > View to filter the dataset in Radiant. The expression should be a string (e.g., "price > 10000")

\item[\code{arr}] Expression to arrange (sort) the data on (e.g., "color, desc(price)")

\item[\code{rows}] Rows to select from the specified dataset

\item[\code{envir}] Environment to extract data from

\item[\code{...}] further arguments passed to or from other methods
\end{ldescription}
\end{Arguments}
%
\begin{Details}
Confusion matrix and additional metrics to evaluate binary classification models. See \url{https://radiant-rstats.github.io/docs/model/evalbin.html} for an example in Radiant
\end{Details}
%
\begin{Value}
A list of results
\end{Value}
%
\begin{SeeAlso}
\code{\LinkA{summary.confusion}{summary.confusion}} to summarize results

\code{\LinkA{plot.confusion}{plot.confusion}} to plot results
\end{SeeAlso}
%
\begin{Examples}
\begin{ExampleCode}
data.frame(buy = dvd$buy, pred1 = runif(20000), pred2 = ifelse(dvd$buy == "yes", 1, 0)) %>%
  confusion(c("pred1", "pred2"), "buy") %>%
  str()
\end{ExampleCode}
\end{Examples}
\HeaderA{crs}{Collaborative Filtering}{crs}
%
\begin{Description}
Collaborative Filtering
\end{Description}
%
\begin{Usage}
\begin{verbatim}
crs(
  dataset,
  id,
  prod,
  pred,
  rate,
  data_filter = "",
  arr = "",
  rows = NULL,
  envir = parent.frame()
)
\end{verbatim}
\end{Usage}
%
\begin{Arguments}
\begin{ldescription}
\item[\code{dataset}] Dataset

\item[\code{id}] String with name of the variable containing user ids

\item[\code{prod}] String with name of the variable with product ids

\item[\code{pred}] Products to predict for

\item[\code{rate}] String with name of the variable with product ratings

\item[\code{data\_filter}] Expression entered in, e.g., Data > View to filter the dataset in Radiant. The expression should be a string (e.g., "training == 1")

\item[\code{arr}] Expression to arrange (sort) the data on (e.g., "color, desc(price)")

\item[\code{rows}] Rows to select from the specified dataset

\item[\code{envir}] Environment to extract data from
\end{ldescription}
\end{Arguments}
%
\begin{Details}
See \url{https://radiant-rstats.github.io/docs/model/crs.html} for an example in Radiant
\end{Details}
%
\begin{Value}
A data.frame with the original data and a new column with predicted ratings
\end{Value}
%
\begin{SeeAlso}
\code{\LinkA{summary.crs}{summary.crs}} to summarize results

\code{\LinkA{plot.crs}{plot.crs}} to plot results if the actual ratings are available
\end{SeeAlso}
%
\begin{Examples}
\begin{ExampleCode}
crs(ratings,
  id = "Users", prod = "Movies", pred = c("M6", "M7", "M8", "M9", "M10"),
  rate = "Ratings", data_filter = "training == 1"
) %>% str()
\end{ExampleCode}
\end{Examples}
\HeaderA{crtree}{Classification and regression trees based on the rpart package}{crtree}
%
\begin{Description}
Classification and regression trees based on the rpart package
\end{Description}
%
\begin{Usage}
\begin{verbatim}
crtree(
  dataset,
  rvar,
  evar,
  type = "",
  lev = "",
  wts = "None",
  minsplit = 2,
  minbucket = round(minsplit/3),
  cp = 0.001,
  pcp = NA,
  nodes = NA,
  K = 10,
  seed = 1234,
  split = "gini",
  prior = NA,
  adjprob = TRUE,
  cost = NA,
  margin = NA,
  check = "",
  data_filter = "",
  arr = "",
  rows = NULL,
  envir = parent.frame()
)
\end{verbatim}
\end{Usage}
%
\begin{Arguments}
\begin{ldescription}
\item[\code{dataset}] Dataset

\item[\code{rvar}] The response variable in the model

\item[\code{evar}] Explanatory variables in the model

\item[\code{type}] Model type (i.e., "classification" or "regression")

\item[\code{lev}] The level in the response variable defined as \_success\_

\item[\code{wts}] Weights to use in estimation

\item[\code{minsplit}] The minimum number of observations that must exist in a node in order for a split to be attempted.

\item[\code{minbucket}] the minimum number of observations in any terminal <leaf> node. If only one of minbucket or minsplit is specified, the code either sets minsplit to minbucket*3 or minbucket to minsplit/3, as appropriate.

\item[\code{cp}] Minimum proportion of root node deviance required for split (default = 0.001)

\item[\code{pcp}] Complexity parameter to use for pruning

\item[\code{nodes}] Maximum size of tree in number of nodes to return

\item[\code{K}] Number of folds use in cross-validation

\item[\code{seed}] Random seed used for cross-validation

\item[\code{split}] Splitting criterion to use (i.e., "gini" or "information")

\item[\code{prior}] Adjust the initial probability for the selected level (e.g., set to .5 in unbalanced samples)

\item[\code{adjprob}] Setting a prior will rescale the predicted probabilities. Set adjprob to TRUE to adjust the probabilities back to their original scale after estimation

\item[\code{cost}] Cost for each treatment (e.g., mailing)

\item[\code{margin}] Margin associated with a successful treatment (e.g., a purchase)

\item[\code{check}] Optional estimation parameters (e.g., "standardize")

\item[\code{data\_filter}] Expression entered in, e.g., Data > View to filter the dataset in Radiant. The expression should be a string (e.g., "price > 10000")

\item[\code{arr}] Expression to arrange (sort) the data on (e.g., "color, desc(price)")

\item[\code{rows}] Rows to select from the specified dataset

\item[\code{envir}] Environment to extract data from
\end{ldescription}
\end{Arguments}
%
\begin{Details}
See \url{https://radiant-rstats.github.io/docs/model/crtree.html} for an example in Radiant
\end{Details}
%
\begin{Value}
A list with all variables defined in crtree as an object of class tree
\end{Value}
%
\begin{SeeAlso}
\code{\LinkA{summary.crtree}{summary.crtree}} to summarize results

\code{\LinkA{plot.crtree}{plot.crtree}} to plot results

\code{\LinkA{predict.crtree}{predict.crtree}} for prediction
\end{SeeAlso}
%
\begin{Examples}
\begin{ExampleCode}
crtree(titanic, "survived", c("pclass", "sex"), lev = "Yes") %>% summary()
result <- crtree(titanic, "survived", c("pclass", "sex")) %>% summary()
result <- crtree(diamonds, "price", c("carat", "clarity"), type = "regression") %>% str()
\end{ExampleCode}
\end{Examples}
\HeaderA{cv.crtree}{Cross-validation for Classification and Regression Trees}{cv.crtree}
%
\begin{Description}
Cross-validation for Classification and Regression Trees
\end{Description}
%
\begin{Usage}
\begin{verbatim}
cv.crtree(
  object,
  K = 5,
  repeats = 1,
  cp,
  pcp = seq(0, 0.01, length.out = 11),
  seed = 1234,
  trace = TRUE,
  fun,
  ...
)
\end{verbatim}
\end{Usage}
%
\begin{Arguments}
\begin{ldescription}
\item[\code{object}] Object of type "rpart" or "crtree" to use as a starting point for cross validation

\item[\code{K}] Number of cross validation passes to use

\item[\code{repeats}] Number of times to repeat the K cross-validation steps

\item[\code{cp}] Complexity parameter used when building the (e.g., 0.0001)

\item[\code{pcp}] Complexity parameter to use for pruning

\item[\code{seed}] Random seed to use as the starting point

\item[\code{trace}] Print progress

\item[\code{fun}] Function to use for model evaluation (e.g., auc for classification or RMSE for regression)

\item[\code{...}] Additional arguments to be passed to 'fun'
\end{ldescription}
\end{Arguments}
%
\begin{Details}
See \url{https://radiant-rstats.github.io/docs/model/crtree.html} for an example in Radiant
\end{Details}
%
\begin{Value}
A data.frame sorted by the mean, sd, min, and max of the performance metric
\end{Value}
%
\begin{SeeAlso}
\code{\LinkA{crtree}{crtree}} to generate an initial model that can be passed to cv.crtree

\code{\LinkA{Rsq}{Rsq}} to calculate an R-squared measure for a regression

\code{\LinkA{RMSE}{RMSE}} to calculate the Root Mean Squared Error for a regression

\code{\LinkA{MAE}{MAE}} to calculate the Mean Absolute Error for a regression

\code{\LinkA{auc}{auc}} to calculate the area under the ROC curve for classification

\code{\LinkA{profit}{profit}} to calculate profits for classification at a cost/margin threshold
\end{SeeAlso}
%
\begin{Examples}
\begin{ExampleCode}
## Not run: 
result <- crtree(dvd, "buy", c("coupon", "purch", "last"))
cv.crtree(result, cp = 0.0001, pcp = seq(0, 0.01, length.out = 11))
cv.crtree(result, cp = 0.0001, pcp = c(0, 0.001, 0.002), fun = profit, cost = 1, margin = 5)
result <- crtree(diamonds, "price", c("carat", "color", "clarity"), type = "regression", cp = 0.001)
cv.crtree(result, cp = 0.001, pcp = seq(0, 0.01, length.out = 11), fun = MAE)

## End(Not run)

\end{ExampleCode}
\end{Examples}
\HeaderA{cv.gbt}{Cross-validation for Gradient Boosted Trees}{cv.gbt}
%
\begin{Description}
Cross-validation for Gradient Boosted Trees
\end{Description}
%
\begin{Usage}
\begin{verbatim}
cv.gbt(
  object,
  K = 5,
  repeats = 1,
  params = list(),
  nrounds = 500,
  early_stopping_rounds = 10,
  nthread = 12,
  train = NULL,
  type = "classification",
  trace = TRUE,
  seed = 1234,
  maximize = NULL,
  fun,
  ...
)
\end{verbatim}
\end{Usage}
%
\begin{Arguments}
\begin{ldescription}
\item[\code{object}] Object of type "gbt" or "ranger"

\item[\code{K}] Number of cross validation passes to use (aka nfold)

\item[\code{repeats}] Repeated cross validation

\item[\code{params}] List of parameters (see XGBoost documentation)

\item[\code{nrounds}] Number of trees to create

\item[\code{early\_stopping\_rounds}] Early stopping rule

\item[\code{nthread}] Number of parallel threads to use. Defaults to 12 if available

\item[\code{train}] An optional xgb.DMatrix object containing the original training data. Not needed when using Radiant's gbt function

\item[\code{type}] Model type ("classification" or "regression")

\item[\code{trace}] Print progress

\item[\code{seed}] Random seed to use as the starting point

\item[\code{maximize}] When a custom function is used, xgb.cv requires the user indicate if the function output should be maximized (TRUE) or minimized (FALSE)

\item[\code{fun}] Function to use for model evaluation (i.e., auc for classification and RMSE for regression)

\item[\code{...}] Additional arguments to be passed to 'fun'
\end{ldescription}
\end{Arguments}
%
\begin{Details}
See \url{https://radiant-rstats.github.io/docs/model/gbt.html} for an example in Radiant
\end{Details}
%
\begin{Value}
A data.frame sorted by the mean of the performance metric
\end{Value}
%
\begin{SeeAlso}
\code{\LinkA{gbt}{gbt}} to generate an initial model that can be passed to cv.gbt

\code{\LinkA{Rsq}{Rsq}} to calculate an R-squared measure for a regression

\code{\LinkA{RMSE}{RMSE}} to calculate the Root Mean Squared Error for a regression

\code{\LinkA{MAE}{MAE}} to calculate the Mean Absolute Error for a regression

\code{\LinkA{auc}{auc}} to calculate the area under the ROC curve for classification

\code{\LinkA{profit}{profit}} to calculate profits for classification at a cost/margin threshold
\end{SeeAlso}
%
\begin{Examples}
\begin{ExampleCode}
## Not run: 
result <- gbt(dvd, "buy", c("coupon", "purch", "last"))
cv.gbt(result, params = list(max_depth = 1:6))
cv.gbt(result, params = list(max_depth = 1:6), fun = "logloss")
cv.gbt(
  result,
  params = list(learning_rate = seq(0.1, 1.0, 0.1)),
  maximize = TRUE, fun = profit, cost = 1, margin = 5
)
result <- gbt(diamonds, "price", c("carat", "color", "clarity"), type = "regression")
cv.gbt(result, params = list(max_depth = 1:2, min_child_weight = 1:2))
cv.gbt(result, params = list(learning_rate = seq(0.1, 0.5, 0.1)), fun = Rsq, maximize = TRUE)
cv.gbt(result, params = list(learning_rate = seq(0.1, 0.5, 0.1)), fun = MAE, maximize = FALSE)

## End(Not run)

\end{ExampleCode}
\end{Examples}
\HeaderA{cv.nn}{Cross-validation for a Neural Network}{cv.nn}
%
\begin{Description}
Cross-validation for a Neural Network
\end{Description}
%
\begin{Usage}
\begin{verbatim}
cv.nn(
  object,
  K = 5,
  repeats = 1,
  decay = seq(0, 1, 0.2),
  size = 1:5,
  seed = 1234,
  trace = TRUE,
  fun,
  ...
)
\end{verbatim}
\end{Usage}
%
\begin{Arguments}
\begin{ldescription}
\item[\code{object}] Object of type "nn" or "nnet"

\item[\code{K}] Number of cross validation passes to use

\item[\code{repeats}] Repeated cross validation

\item[\code{decay}] Parameter decay

\item[\code{size}] Number of units (nodes) in the hidden layer

\item[\code{seed}] Random seed to use as the starting point

\item[\code{trace}] Print progress

\item[\code{fun}] Function to use for model evaluation (i.e., auc for classification and RMSE for regression)

\item[\code{...}] Additional arguments to be passed to 'fun'
\end{ldescription}
\end{Arguments}
%
\begin{Details}
See \url{https://radiant-rstats.github.io/docs/model/nn.html} for an example in Radiant
\end{Details}
%
\begin{Value}
A data.frame sorted by the mean of the performance metric
\end{Value}
%
\begin{SeeAlso}
\code{\LinkA{nn}{nn}} to generate an initial model that can be passed to cv.nn

\code{\LinkA{Rsq}{Rsq}} to calculate an R-squared measure for a regression

\code{\LinkA{RMSE}{RMSE}} to calculate the Root Mean Squared Error for a regression

\code{\LinkA{MAE}{MAE}} to calculate the Mean Absolute Error for a regression

\code{\LinkA{auc}{auc}} to calculate the area under the ROC curve for classification

\code{\LinkA{profit}{profit}} to calculate profits for classification at a cost/margin threshold
\end{SeeAlso}
%
\begin{Examples}
\begin{ExampleCode}
## Not run: 
result <- nn(dvd, "buy", c("coupon", "purch", "last"))
cv.nn(result, decay = seq(0, 1, .5), size = 1:2)
cv.nn(result, decay = seq(0, 1, .5), size = 1:2, fun = profit, cost = 1, margin = 5)
result <- nn(diamonds, "price", c("carat", "color", "clarity"), type = "regression")
cv.nn(result, decay = seq(0, 1, .5), size = 1:2)
cv.nn(result, decay = seq(0, 1, .5), size = 1:2, fun = Rsq)

## End(Not run)

\end{ExampleCode}
\end{Examples}
\HeaderA{cv.rforest}{Cross-validation for a Random Forest}{cv.rforest}
%
\begin{Description}
Cross-validation for a Random Forest
\end{Description}
%
\begin{Usage}
\begin{verbatim}
cv.rforest(
  object,
  K = 5,
  repeats = 1,
  mtry = 1:5,
  num.trees = NULL,
  min.node.size = 1,
  sample.fraction = NA,
  trace = TRUE,
  seed = 1234,
  fun,
  ...
)
\end{verbatim}
\end{Usage}
%
\begin{Arguments}
\begin{ldescription}
\item[\code{object}] Object of type "rforest" or "ranger"

\item[\code{K}] Number of cross validation passes to use

\item[\code{repeats}] Repeated cross validation

\item[\code{mtry}] Number of variables to possibly split at in each node. Default is the (rounded down) square root of the number variables

\item[\code{num.trees}] Number of trees to create

\item[\code{min.node.size}] Minimal node size

\item[\code{sample.fraction}] Fraction of observations to sample. Default is 1 for sampling with replacement and 0.632 for sampling without replacement

\item[\code{trace}] Print progress

\item[\code{seed}] Random seed to use as the starting point

\item[\code{fun}] Function to use for model evaluation (i.e., auc for classification and RMSE for regression)

\item[\code{...}] Additional arguments to be passed to 'fun'
\end{ldescription}
\end{Arguments}
%
\begin{Details}
See \url{https://radiant-rstats.github.io/docs/model/rforest.html} for an example in Radiant
\end{Details}
%
\begin{Value}
A data.frame sorted by the mean of the performance metric
\end{Value}
%
\begin{SeeAlso}
\code{\LinkA{rforest}{rforest}} to generate an initial model that can be passed to cv.rforest

\code{\LinkA{Rsq}{Rsq}} to calculate an R-squared measure for a regression

\code{\LinkA{RMSE}{RMSE}} to calculate the Root Mean Squared Error for a regression

\code{\LinkA{MAE}{MAE}} to calculate the Mean Absolute Error for a regression

\code{\LinkA{auc}{auc}} to calculate the area under the ROC curve for classification

\code{\LinkA{profit}{profit}} to calculate profits for classification at a cost/margin threshold
\end{SeeAlso}
%
\begin{Examples}
\begin{ExampleCode}
## Not run: 
result <- rforest(dvd, "buy", c("coupon", "purch", "last"))
cv.rforest(
  result,
  mtry = 1:3, min.node.size = seq(1, 10, 5),
  num.trees = c(100, 200), sample.fraction = 0.632
)
result <- rforest(titanic, "survived", c("pclass", "sex"), max.depth = 1)
cv.rforest(result, mtry = 1:3, min.node.size = seq(1, 10, 5))
cv.rforest(result, mtry = 1:3, num.trees = c(100, 200), fun = profit, cost = 1, margin = 5)
result <- rforest(diamonds, "price", c("carat", "color", "clarity"), type = "regression")
cv.rforest(result, mtry = 1:3, min.node.size = 1)
cv.rforest(result, mtry = 1:3, min.node.size = 1, fun = Rsq)

## End(Not run)

\end{ExampleCode}
\end{Examples}
\HeaderA{direct\_marketing}{Direct marketing data}{direct.Rul.marketing}
\keyword{datasets}{direct\_marketing}
%
\begin{Description}
Direct marketing data
\end{Description}
%
\begin{Usage}
\begin{verbatim}
data(direct_marketing)
\end{verbatim}
\end{Usage}
%
\begin{Format}
A data frame with 1,000 rows and 12 variables
\end{Format}
%
\begin{Details}
Description provided in attr(direct\_marketing, "description")
\end{Details}
\HeaderA{dtree}{Create a decision tree}{dtree}
%
\begin{Description}
Create a decision tree
\end{Description}
%
\begin{Usage}
\begin{verbatim}
dtree(yl, opt = "max", base = character(0), envir = parent.frame())
\end{verbatim}
\end{Usage}
%
\begin{Arguments}
\begin{ldescription}
\item[\code{yl}] A yaml string or a list (e.g., from yaml::yaml.load\_file())

\item[\code{opt}] Find the maximum ("max") or minimum ("min") value for each decision node

\item[\code{base}] List of variable definitions from a base tree used when calling a sub-tree

\item[\code{envir}] Environment to extract data from
\end{ldescription}
\end{Arguments}
%
\begin{Details}
See \url{https://radiant-rstats.github.io/docs/model/dtree.html} for an example in Radiant
\end{Details}
%
\begin{Value}
A list with the initial tree, the calculated tree, and a data.frame with results (i.e., payoffs, probabilities, etc.)
\end{Value}
%
\begin{SeeAlso}
\code{\LinkA{summary.dtree}{summary.dtree}} to summarize results

\code{\LinkA{plot.dtree}{plot.dtree}} to plot results

\code{\LinkA{sensitivity.dtree}{sensitivity.dtree}} to plot results
\end{SeeAlso}
%
\begin{Examples}
\begin{ExampleCode}
yaml::as.yaml(movie_contract) %>% cat()
dtree(movie_contract, opt = "max") %>% summary(output = TRUE)
dtree(movie_contract)$payoff
dtree(movie_contract)$prob
dtree(movie_contract)$solution_df

\end{ExampleCode}
\end{Examples}
\HeaderA{dtree\_parser}{Parse yaml input for dtree to provide (more) useful error messages}{dtree.Rul.parser}
%
\begin{Description}
Parse yaml input for dtree to provide (more) useful error messages
\end{Description}
%
\begin{Usage}
\begin{verbatim}
dtree_parser(yl)
\end{verbatim}
\end{Usage}
%
\begin{Arguments}
\begin{ldescription}
\item[\code{yl}] A yaml string
\end{ldescription}
\end{Arguments}
%
\begin{Details}
See \url{https://radiant-rstats.github.io/docs/model/dtree.html} for an example in Radiant
\end{Details}
%
\begin{Value}
An updated yaml string or a vector messages to return to the users
\end{Value}
%
\begin{SeeAlso}
\code{\LinkA{dtree}{dtree}} to calculate tree

\code{\LinkA{summary.dtree}{summary.dtree}} to summarize results

\code{\LinkA{plot.dtree}{plot.dtree}} to plot results
\end{SeeAlso}
\HeaderA{dvd}{Data on DVD sales}{dvd}
\keyword{datasets}{dvd}
%
\begin{Description}
Data on DVD sales
\end{Description}
%
\begin{Usage}
\begin{verbatim}
data(dvd)
\end{verbatim}
\end{Usage}
%
\begin{Format}
A data frame with 20,000 rows and 4 variables
\end{Format}
%
\begin{Details}
Binary purchase response to coupon value. Description provided in attr(dvd,"description")
\end{Details}
\HeaderA{evalbin}{Evaluate the performance of different (binary) classification models}{evalbin}
%
\begin{Description}
Evaluate the performance of different (binary) classification models
\end{Description}
%
\begin{Usage}
\begin{verbatim}
evalbin(
  dataset,
  pred,
  rvar,
  lev = "",
  qnt = 10,
  cost = 1,
  margin = 2,
  scale = 1,
  train = "All",
  data_filter = "",
  arr = "",
  rows = NULL,
  envir = parent.frame()
)
\end{verbatim}
\end{Usage}
%
\begin{Arguments}
\begin{ldescription}
\item[\code{dataset}] Dataset

\item[\code{pred}] Predictions or predictors

\item[\code{rvar}] Response variable

\item[\code{lev}] The level in the response variable defined as success

\item[\code{qnt}] Number of bins to create

\item[\code{cost}] Cost for each connection (e.g., email or mailing)

\item[\code{margin}] Margin on each customer purchase

\item[\code{scale}] Scaling factor to apply to calculations

\item[\code{train}] Use data from training ("Training"), test ("Test"), both ("Both"), or all data ("All") to evaluate model evalbin

\item[\code{data\_filter}] Expression entered in, e.g., Data > View to filter the dataset in Radiant. The expression should be a string (e.g., "price > 10000")

\item[\code{arr}] Expression to arrange (sort) the data on (e.g., "color, desc(price)")

\item[\code{rows}] Rows to select from the specified dataset

\item[\code{envir}] Environment to extract data from
\end{ldescription}
\end{Arguments}
%
\begin{Details}
Evaluate different (binary) classification models based on predictions. See \url{https://radiant-rstats.github.io/docs/model/evalbin.html} for an example in Radiant
\end{Details}
%
\begin{Value}
A list of results
\end{Value}
%
\begin{SeeAlso}
\code{\LinkA{summary.evalbin}{summary.evalbin}} to summarize results

\code{\LinkA{plot.evalbin}{plot.evalbin}} to plot results
\end{SeeAlso}
%
\begin{Examples}
\begin{ExampleCode}
data.frame(buy = dvd$buy, pred1 = runif(20000), pred2 = ifelse(dvd$buy == "yes", 1, 0)) %>%
  evalbin(c("pred1", "pred2"), "buy") %>%
  str()
\end{ExampleCode}
\end{Examples}
\HeaderA{evalreg}{Evaluate the performance of different regression models}{evalreg}
%
\begin{Description}
Evaluate the performance of different regression models
\end{Description}
%
\begin{Usage}
\begin{verbatim}
evalreg(
  dataset,
  pred,
  rvar,
  train = "All",
  data_filter = "",
  arr = "",
  rows = NULL,
  envir = parent.frame()
)
\end{verbatim}
\end{Usage}
%
\begin{Arguments}
\begin{ldescription}
\item[\code{dataset}] Dataset

\item[\code{pred}] Predictions or predictors

\item[\code{rvar}] Response variable

\item[\code{train}] Use data from training ("Training"), test ("Test"), both ("Both"), or all data ("All") to evaluate model evalreg

\item[\code{data\_filter}] Expression entered in, e.g., Data > View to filter the dataset in Radiant. The expression should be a string (e.g., "training == 1")

\item[\code{arr}] Expression to arrange (sort) the data on (e.g., "color, desc(price)")

\item[\code{rows}] Rows to select from the specified dataset

\item[\code{envir}] Environment to extract data from
\end{ldescription}
\end{Arguments}
%
\begin{Details}
Evaluate different regression models based on predictions. See \url{https://radiant-rstats.github.io/docs/model/evalreg.html} for an example in Radiant
\end{Details}
%
\begin{Value}
A list of results
\end{Value}
%
\begin{SeeAlso}
\code{\LinkA{summary.evalreg}{summary.evalreg}} to summarize results

\code{\LinkA{plot.evalreg}{plot.evalreg}} to plot results
\end{SeeAlso}
%
\begin{Examples}
\begin{ExampleCode}
data.frame(price = diamonds$price, pred1 = rnorm(3000), pred2 = diamonds$price) %>%
  evalreg(pred = c("pred1", "pred2"), "price") %>%
  str()

\end{ExampleCode}
\end{Examples}
\HeaderA{find\_max}{Find maximum value of a vector}{find.Rul.max}
%
\begin{Description}
Find maximum value of a vector
\end{Description}
%
\begin{Usage}
\begin{verbatim}
find_max(x, y)
\end{verbatim}
\end{Usage}
%
\begin{Arguments}
\begin{ldescription}
\item[\code{x}] Variable to find the maximum for

\item[\code{y}] Variable to find the value for at the maximum of var
\end{ldescription}
\end{Arguments}
%
\begin{Details}
Find the value of y at the maximum value of x
\end{Details}
%
\begin{Value}
Value of val at the maximum of var
\end{Value}
%
\begin{Examples}
\begin{ExampleCode}
find_max(1:10, 21:30)

\end{ExampleCode}
\end{Examples}
\HeaderA{find\_min}{Find minimum value of a vector}{find.Rul.min}
%
\begin{Description}
Find minimum value of a vector
\end{Description}
%
\begin{Usage}
\begin{verbatim}
find_min(x, y)
\end{verbatim}
\end{Usage}
%
\begin{Arguments}
\begin{ldescription}
\item[\code{x}] Variable to find the minimum for

\item[\code{y}] Variable to find the value for at the maximum of var
\end{ldescription}
\end{Arguments}
%
\begin{Details}
Find the value of y at the minimum value of x
\end{Details}
%
\begin{Value}
Value of val at the minimum of var
\end{Value}
%
\begin{Examples}
\begin{ExampleCode}
find_min(1:10, 21:30)

\end{ExampleCode}
\end{Examples}
\HeaderA{gbt}{Gradient Boosted Trees using XGBoost}{gbt}
%
\begin{Description}
Gradient Boosted Trees using XGBoost
\end{Description}
%
\begin{Usage}
\begin{verbatim}
gbt(
  dataset,
  rvar,
  evar,
  type = "classification",
  lev = "",
  max_depth = 6,
  learning_rate = 0.3,
  min_split_loss = 0,
  min_child_weight = 1,
  subsample = 1,
  nrounds = 100,
  early_stopping_rounds = 10,
  nthread = 12,
  wts = "None",
  seed = NA,
  data_filter = "",
  arr = "",
  rows = NULL,
  envir = parent.frame(),
  ...
)
\end{verbatim}
\end{Usage}
%
\begin{Arguments}
\begin{ldescription}
\item[\code{dataset}] Dataset

\item[\code{rvar}] The response variable in the model

\item[\code{evar}] Explanatory variables in the model

\item[\code{type}] Model type (i.e., "classification" or "regression")

\item[\code{lev}] Level to use as the first column in prediction output

\item[\code{max\_depth}] Maximum 'depth' of tree

\item[\code{learning\_rate}] Learning rate (eta)

\item[\code{min\_split\_loss}] Minimal improvement (gamma)

\item[\code{min\_child\_weight}] Minimum number of instances allowed in each node

\item[\code{subsample}] Subsample ratio of the training instances (0-1)

\item[\code{nrounds}] Number of trees to create

\item[\code{early\_stopping\_rounds}] Early stopping rule

\item[\code{nthread}] Number of parallel threads to use. Defaults to 12 if available

\item[\code{wts}] Weights to use in estimation

\item[\code{seed}] Random seed to use as the starting point

\item[\code{data\_filter}] Expression entered in, e.g., Data > View to filter the dataset in Radiant. The expression should be a string (e.g., "price > 10000")

\item[\code{arr}] Expression to arrange (sort) the data on (e.g., "color, desc(price)")

\item[\code{rows}] Rows to select from the specified dataset

\item[\code{envir}] Environment to extract data from

\item[\code{...}] Further arguments to pass to xgboost
\end{ldescription}
\end{Arguments}
%
\begin{Details}
See \url{https://radiant-rstats.github.io/docs/model/gbt.html} for an example in Radiant
\end{Details}
%
\begin{Value}
A list with all variables defined in gbt as an object of class gbt
\end{Value}
%
\begin{SeeAlso}
\code{\LinkA{summary.gbt}{summary.gbt}} to summarize results

\code{\LinkA{plot.gbt}{plot.gbt}} to plot results

\code{\LinkA{predict.gbt}{predict.gbt}} for prediction
\end{SeeAlso}
%
\begin{Examples}
\begin{ExampleCode}
## Not run: 
gbt(titanic, "survived", c("pclass", "sex"), lev = "Yes") %>% summary()
gbt(titanic, "survived", c("pclass", "sex")) %>% str()

## End(Not run)
gbt(
  titanic, "survived", c("pclass", "sex"), lev = "Yes",
  early_stopping_rounds = 0, nthread = 1
) %>% summary()
gbt(
  titanic, "survived", c("pclass", "sex"),
  early_stopping_rounds = 0, nthread = 1
) %>% str()
gbt(
  titanic, "survived", c("pclass", "sex"),
  eval_metric = paste0("error@", 0.5 / 6), nthread = 1
) %>% str()
gbt(
  diamonds, "price", c("carat", "clarity"), type = "regression", nthread = 1
) %>% summary()

\end{ExampleCode}
\end{Examples}
\HeaderA{houseprices}{Houseprices}{houseprices}
\keyword{datasets}{houseprices}
%
\begin{Description}
Houseprices
\end{Description}
%
\begin{Usage}
\begin{verbatim}
data(houseprices)
\end{verbatim}
\end{Usage}
%
\begin{Format}
A data frame with 128 home sales and 6 variables
\end{Format}
%
\begin{Details}
Description provided in attr(houseprices, "description")
\end{Details}
\HeaderA{ideal}{Ideal data for linear regression}{ideal}
\keyword{datasets}{ideal}
%
\begin{Description}
Ideal data for linear regression
\end{Description}
%
\begin{Usage}
\begin{verbatim}
data(ideal)
\end{verbatim}
\end{Usage}
%
\begin{Format}
A data frame with 1,000 rows and 4 variables
\end{Format}
%
\begin{Details}
Description provided in attr(ideal,  "description")
\end{Details}
\HeaderA{kaggle\_uplift}{Kaggle uplift}{kaggle.Rul.uplift}
\keyword{datasets}{kaggle\_uplift}
%
\begin{Description}
Kaggle uplift
\end{Description}
%
\begin{Usage}
\begin{verbatim}
data(kaggle_uplift)
\end{verbatim}
\end{Usage}
%
\begin{Format}
A data frame with 1,000 rows and 22 variables
\end{Format}
%
\begin{Details}
Use uplift modeling to quantify the effectiveness of an experimental treatment
\end{Details}
\HeaderA{ketchup}{Data on ketchup choices}{ketchup}
\keyword{datasets}{ketchup}
%
\begin{Description}
Data on ketchup choices
\end{Description}
%
\begin{Usage}
\begin{verbatim}
data(ketchup)
\end{verbatim}
\end{Usage}
%
\begin{Format}
A data frame with 2,798 rows and 14 variables
\end{Format}
%
\begin{Details}
Choice behavior for a sample of 300 individuals in a panel of households in Springfield, Missouri (USA). Description provided in attr(ketchup,"description")
\end{Details}
\HeaderA{logistic}{Logistic regression}{logistic}
%
\begin{Description}
Logistic regression
\end{Description}
%
\begin{Usage}
\begin{verbatim}
logistic(
  dataset,
  rvar,
  evar,
  lev = "",
  int = "",
  wts = "None",
  check = "",
  form,
  ci_type,
  data_filter = "",
  arr = "",
  rows = NULL,
  envir = parent.frame()
)
\end{verbatim}
\end{Usage}
%
\begin{Arguments}
\begin{ldescription}
\item[\code{dataset}] Dataset

\item[\code{rvar}] The response variable in the model

\item[\code{evar}] Explanatory variables in the model

\item[\code{lev}] The level in the response variable defined as \_success\_

\item[\code{int}] Interaction term to include in the model

\item[\code{wts}] Weights to use in estimation

\item[\code{check}] Use "standardize" to see standardized coefficient estimates. Use "stepwise-backward" (or "stepwise-forward", or "stepwise-both") to apply step-wise selection of variables in estimation. Add "robust" for robust estimation of standard errors (HC1)

\item[\code{form}] Optional formula to use instead of rvar, evar, and int

\item[\code{ci\_type}] To use the profile-likelihood (rather than Wald) for confidence intervals use "profile". For datasets with more than 5,000 rows the Wald method will be used, unless "profile" is explicitly set

\item[\code{data\_filter}] Expression entered in, e.g., Data > View to filter the dataset in Radiant. The expression should be a string (e.g., "price > 10000")

\item[\code{arr}] Expression to arrange (sort) the data on (e.g., "color, desc(price)")

\item[\code{rows}] Rows to select from the specified dataset

\item[\code{envir}] Environment to extract data from
\end{ldescription}
\end{Arguments}
%
\begin{Details}
See \url{https://radiant-rstats.github.io/docs/model/logistic.html} for an example in Radiant
\end{Details}
%
\begin{Value}
A list with all variables defined in logistic as an object of class logistic
\end{Value}
%
\begin{SeeAlso}
\code{\LinkA{summary.logistic}{summary.logistic}} to summarize the results

\code{\LinkA{plot.logistic}{plot.logistic}} to plot the results

\code{\LinkA{predict.logistic}{predict.logistic}} to generate predictions

\code{\LinkA{plot.model.predict}{plot.model.predict}} to plot prediction output
\end{SeeAlso}
%
\begin{Examples}
\begin{ExampleCode}
logistic(titanic, "survived", c("pclass", "sex"), lev = "Yes") %>% summary()
logistic(titanic, "survived", c("pclass", "sex")) %>% str()
\end{ExampleCode}
\end{Examples}
\HeaderA{MAE}{Mean Absolute Error}{MAE}
%
\begin{Description}
Mean Absolute Error
\end{Description}
%
\begin{Usage}
\begin{verbatim}
MAE(pred, rvar)
\end{verbatim}
\end{Usage}
%
\begin{Arguments}
\begin{ldescription}
\item[\code{pred}] Prediction (vector)

\item[\code{rvar}] Response (vector)
\end{ldescription}
\end{Arguments}
%
\begin{Value}
Mean Absolute Error
\end{Value}
\HeaderA{minmax}{Calculate min and max before standardization}{minmax}
%
\begin{Description}
Calculate min and max before standardization
\end{Description}
%
\begin{Usage}
\begin{verbatim}
minmax(dataset)
\end{verbatim}
\end{Usage}
%
\begin{Arguments}
\begin{ldescription}
\item[\code{dataset}] Data frame
\end{ldescription}
\end{Arguments}
%
\begin{Value}
Data frame min and max attributes
\end{Value}
\HeaderA{mnl}{Multinomial logistic regression}{mnl}
%
\begin{Description}
Multinomial logistic regression
\end{Description}
%
\begin{Usage}
\begin{verbatim}
mnl(
  dataset,
  rvar,
  evar,
  lev = "",
  int = "",
  wts = "None",
  check = "",
  data_filter = "",
  arr = "",
  rows = NULL,
  envir = parent.frame()
)
\end{verbatim}
\end{Usage}
%
\begin{Arguments}
\begin{ldescription}
\item[\code{dataset}] Dataset

\item[\code{rvar}] The response variable in the model

\item[\code{evar}] Explanatory variables in the model

\item[\code{lev}] The level in the response variable to use as the baseline

\item[\code{int}] Interaction term to include in the model

\item[\code{wts}] Weights to use in estimation

\item[\code{check}] Use "standardize" to see standardized coefficient estimates. Use "stepwise-backward" (or "stepwise-forward", or "stepwise-both") to apply step-wise selection of variables in estimation.

\item[\code{data\_filter}] Expression entered in, e.g., Data > View to filter the dataset in Radiant. The expression should be a string (e.g., "price > 10000")

\item[\code{arr}] Expression to arrange (sort) the data on (e.g., "color, desc(price)")

\item[\code{rows}] Rows to select from the specified dataset

\item[\code{envir}] Environment to extract data from
\end{ldescription}
\end{Arguments}
%
\begin{Details}
See \url{https://radiant-rstats.github.io/docs/model/mnl.html} for an example in Radiant
\end{Details}
%
\begin{Value}
A list with all variables defined in mnl as an object of class mnl
\end{Value}
%
\begin{SeeAlso}
\code{\LinkA{summary.mnl}{summary.mnl}} to summarize the results

\code{\LinkA{plot.mnl}{plot.mnl}} to plot the results

\code{\LinkA{predict.mnl}{predict.mnl}} to generate predictions

\code{\LinkA{plot.model.predict}{plot.model.predict}} to plot prediction output
\end{SeeAlso}
%
\begin{Examples}
\begin{ExampleCode}
result <- mnl(
  ketchup,
  rvar = "choice",
  evar = c("price.heinz28", "price.heinz32", "price.heinz41", "price.hunts32"),
  lev = "heinz28"
)
str(result)

\end{ExampleCode}
\end{Examples}
\HeaderA{movie\_contract}{Movie contract decision tree}{movie.Rul.contract}
\keyword{datasets}{movie\_contract}
%
\begin{Description}
Movie contract decision tree
\end{Description}
%
\begin{Usage}
\begin{verbatim}
data(movie_contract)
\end{verbatim}
\end{Usage}
%
\begin{Format}
A nested list for decision and chance nodes, probabilities and payoffs
\end{Format}
%
\begin{Details}
Use decision analysis to create a decision tree for an actor facing a contract decision
\end{Details}
\HeaderA{nb}{Naive Bayes using e1071::naiveBayes}{nb}
%
\begin{Description}
Naive Bayes using e1071::naiveBayes
\end{Description}
%
\begin{Usage}
\begin{verbatim}
nb(
  dataset,
  rvar,
  evar,
  laplace = 0,
  data_filter = "",
  arr = "",
  rows = NULL,
  envir = parent.frame()
)
\end{verbatim}
\end{Usage}
%
\begin{Arguments}
\begin{ldescription}
\item[\code{dataset}] Dataset

\item[\code{rvar}] The response variable in the logit (probit) model

\item[\code{evar}] Explanatory variables in the model

\item[\code{laplace}] Positive double controlling Laplace smoothing. The default (0) disables Laplace smoothing.

\item[\code{data\_filter}] Expression entered in, e.g., Data > View to filter the dataset in Radiant. The expression should be a string (e.g., "price > 10000")

\item[\code{arr}] Expression to arrange (sort) the data on (e.g., "color, desc(price)")

\item[\code{rows}] Rows to select from the specified dataset

\item[\code{envir}] Environment to extract data from
\end{ldescription}
\end{Arguments}
%
\begin{Details}
See \url{https://radiant-rstats.github.io/docs/model/nb.html} for an example in Radiant
\end{Details}
%
\begin{Value}
A list with all variables defined in nb as an object of class nb
\end{Value}
%
\begin{SeeAlso}
\code{\LinkA{summary.nb}{summary.nb}} to summarize results

\code{\LinkA{plot.nb}{plot.nb}} to plot results

\code{\LinkA{predict.nb}{predict.nb}} for prediction
\end{SeeAlso}
%
\begin{Examples}
\begin{ExampleCode}
nb(titanic, "survived", c("pclass", "sex", "age")) %>% summary()
nb(titanic, "survived", c("pclass", "sex", "age")) %>% str()

\end{ExampleCode}
\end{Examples}
\HeaderA{nn}{Neural Networks using nnet}{nn}
%
\begin{Description}
Neural Networks using nnet
\end{Description}
%
\begin{Usage}
\begin{verbatim}
nn(
  dataset,
  rvar,
  evar,
  type = "classification",
  lev = "",
  size = 1,
  decay = 0.5,
  wts = "None",
  seed = NA,
  check = "standardize",
  form,
  data_filter = "",
  arr = "",
  rows = NULL,
  envir = parent.frame()
)
\end{verbatim}
\end{Usage}
%
\begin{Arguments}
\begin{ldescription}
\item[\code{dataset}] Dataset

\item[\code{rvar}] The response variable in the model

\item[\code{evar}] Explanatory variables in the model

\item[\code{type}] Model type (i.e., "classification" or "regression")

\item[\code{lev}] The level in the response variable defined as \_success\_

\item[\code{size}] Number of units (nodes) in the hidden layer

\item[\code{decay}] Parameter decay

\item[\code{wts}] Weights to use in estimation

\item[\code{seed}] Random seed to use as the starting point

\item[\code{check}] Optional estimation parameters ("standardize" is the default)

\item[\code{form}] Optional formula to use instead of rvar and evar

\item[\code{data\_filter}] Expression entered in, e.g., Data > View to filter the dataset in Radiant. The expression should be a string (e.g., "price > 10000")

\item[\code{arr}] Expression to arrange (sort) the data on (e.g., "color, desc(price)")

\item[\code{rows}] Rows to select from the specified dataset

\item[\code{envir}] Environment to extract data from
\end{ldescription}
\end{Arguments}
%
\begin{Details}
See \url{https://radiant-rstats.github.io/docs/model/nn.html} for an example in Radiant
\end{Details}
%
\begin{Value}
A list with all variables defined in nn as an object of class nn
\end{Value}
%
\begin{SeeAlso}
\code{\LinkA{summary.nn}{summary.nn}} to summarize results

\code{\LinkA{plot.nn}{plot.nn}} to plot results

\code{\LinkA{predict.nn}{predict.nn}} for prediction
\end{SeeAlso}
%
\begin{Examples}
\begin{ExampleCode}
nn(titanic, "survived", c("pclass", "sex"), lev = "Yes") %>% summary()
nn(titanic, "survived", c("pclass", "sex")) %>% str()
nn(diamonds, "price", c("carat", "clarity"), type = "regression") %>% summary()
\end{ExampleCode}
\end{Examples}
\HeaderA{onehot}{One hot encoding of data.frames}{onehot}
%
\begin{Description}
One hot encoding of data.frames
\end{Description}
%
\begin{Usage}
\begin{verbatim}
onehot(dataset, all = FALSE, df = FALSE)
\end{verbatim}
\end{Usage}
%
\begin{Arguments}
\begin{ldescription}
\item[\code{dataset}] Dataset to endcode

\item[\code{all}] Extract all factor levels (e.g., for tree-based models)

\item[\code{df}] Return a data.frame (tibble)
\end{ldescription}
\end{Arguments}
%
\begin{Examples}
\begin{ExampleCode}
head(onehot(diamonds, df = TRUE))
head(onehot(diamonds, all = TRUE, df = TRUE))
\end{ExampleCode}
\end{Examples}
\HeaderA{pdp\_plot}{Create Partial Dependence Plots}{pdp.Rul.plot}
%
\begin{Description}
Create Partial Dependence Plots
\end{Description}
%
\begin{Usage}
\begin{verbatim}
pdp_plot(
  x,
  plot_list = list(),
  incl,
  incl_int,
  fix = TRUE,
  hline = TRUE,
  nr = 20,
  pdp_range = c(0.025, 0.975),
  minq = NULL,
  maxq = NULL
)
\end{verbatim}
\end{Usage}
%
\begin{Arguments}
\begin{ldescription}
\item[\code{x}] Return value from a model

\item[\code{plot\_list}] List used to store plots

\item[\code{incl}] Which variables to include in PDP plots

\item[\code{incl\_int}] Which interactions to investigate in PDP plots

\item[\code{fix}] Set the desired limited on yhat or have it calculated automatically.
Set to FALSE to have y-axis limits set by ggplot2 for each plot

\item[\code{hline}] Add a dashed horizontal line at the mean response. Set to FALSE to suppress,
or a numeric value to draw the line at that specific value

\item[\code{nr}] Number of values to use to generate predictions for a numeric explanatory variable

\item[\code{pdp\_range}] Numeric vector \code{c(lo, hi)} giving the percentile range used to trim the
x-axis for numeric predictors (default \code{c(0.025, 0.975)})

\item[\code{minq}] Deprecated. Use \code{pdp\_range[1]} instead

\item[\code{maxq}] Deprecated. Use \code{pdp\_range[2]} instead
\end{ldescription}
\end{Arguments}
\HeaderA{plot.confusion}{Plot method for the confusion matrix}{plot.confusion}
%
\begin{Description}
Plot method for the confusion matrix
\end{Description}
%
\begin{Usage}
\begin{verbatim}
## S3 method for class 'confusion'
plot(
  x,
  vars = c("kappa", "index", "ROME", "AUC"),
  scale_y = TRUE,
  size = 13,
  ...
)
\end{verbatim}
\end{Usage}
%
\begin{Arguments}
\begin{ldescription}
\item[\code{x}] Return value from \code{\LinkA{confusion}{confusion}}

\item[\code{vars}] Measures to plot, i.e., one or more of "TP", "FP", "TN", "FN", "total", "TPR", "TNR", "precision", "accuracy", "kappa", "profit", "index", "ROME", "contact", "AUC"

\item[\code{scale\_y}] Free scale in faceted plot of the confusion matrix (TRUE or FALSE)

\item[\code{size}] Font size used

\item[\code{...}] further arguments passed to or from other methods
\end{ldescription}
\end{Arguments}
%
\begin{Details}
See \url{https://radiant-rstats.github.io/docs/model/evalbin.html} for an example in Radiant
\end{Details}
%
\begin{SeeAlso}
\code{\LinkA{confusion}{confusion}} to generate results

\code{\LinkA{summary.confusion}{summary.confusion}} to summarize results
\end{SeeAlso}
%
\begin{Examples}
\begin{ExampleCode}
data.frame(buy = dvd$buy, pred1 = runif(20000), pred2 = ifelse(dvd$buy == "yes", 1, 0)) %>%
  confusion(c("pred1", "pred2"), "buy") %>%
  plot()
\end{ExampleCode}
\end{Examples}
\HeaderA{plot.crs}{Plot method for the crs function}{plot.crs}
%
\begin{Description}
Plot method for the crs function
\end{Description}
%
\begin{Usage}
\begin{verbatim}
## S3 method for class 'crs'
plot(x, ...)
\end{verbatim}
\end{Usage}
%
\begin{Arguments}
\begin{ldescription}
\item[\code{x}] Return value from \code{\LinkA{crs}{crs}}

\item[\code{...}] further arguments passed to or from other methods
\end{ldescription}
\end{Arguments}
%
\begin{Details}
Plot that compares actual to predicted ratings. See \url{https://radiant-rstats.github.io/docs/model/crs.html} for an example in Radiant
\end{Details}
%
\begin{SeeAlso}
\code{\LinkA{crs}{crs}} to generate results

\code{\LinkA{summary.crs}{summary.crs}} to summarize results
\end{SeeAlso}
\HeaderA{plot.crtree}{Plot method for the crtree function}{plot.crtree}
%
\begin{Description}
Plot method for the crtree function
\end{Description}
%
\begin{Usage}
\begin{verbatim}
## S3 method for class 'crtree'
plot(
  x,
  plots = "tree",
  orient = "LR",
  width = "900px",
  labs = TRUE,
  nrobs = Inf,
  dec = 2,
  incl = NULL,
  incl_int = NULL,
  hline = TRUE,
  pdp_range = c(0.025, 0.975),
  minq = NULL,
  maxq = NULL,
  shiny = FALSE,
  custom = FALSE,
  ...
)
\end{verbatim}
\end{Usage}
%
\begin{Arguments}
\begin{ldescription}
\item[\code{x}] Return value from \code{\LinkA{crtree}{crtree}}

\item[\code{plots}] Plots to produce for the specified rpart tree. "tree" shows a tree diagram. "prune" shows a line graph to evaluate appropriate tree pruning. "imp" shows a variable importance plot

\item[\code{orient}] Plot orientation for tree: LR for vertical and TD for horizontal

\item[\code{width}] Plot width in pixels for tree (default is "900px")

\item[\code{labs}] Use factor labels in plot (TRUE) or revert to default letters used by tree (FALSE)

\item[\code{nrobs}] Number of data points to show in dashboard scatter plots (-1 for all)

\item[\code{dec}] Decimal places to round results to

\item[\code{incl}] Which variables to include in a coefficient plot or PDP plot

\item[\code{incl\_int}] Which interactions to investigate in PDP plots

\item[\code{hline}] Add a dashed horizontal line at the mean response (TRUE/FALSE or a numeric value)

\item[\code{pdp\_range}] Numeric vector \code{c(lo, hi)} giving the percentile range for PDP/Prediction plot x-axes (default \code{c(0.025, 0.975)})

\item[\code{minq}] Deprecated. Use \code{pdp\_range[1]} instead

\item[\code{maxq}] Deprecated. Use \code{pdp\_range[2]} instead

\item[\code{shiny}] Did the function call originate inside a shiny app

\item[\code{custom}] Logical (TRUE, FALSE) to indicate if ggplot object (or list of ggplot objects) should be returned. This option can be used to customize plots (e.g., add a title, change x and y labels, etc.). See examples and \url{https://ggplot2.tidyverse.org} for options.

\item[\code{...}] further arguments passed to or from other methods
\end{ldescription}
\end{Arguments}
%
\begin{Details}
Plot a decision tree using mermaid, permutation plots , prediction plots, or partial dependence plots. For regression trees, a residual dashboard can be plotted. See \url{https://radiant-rstats.github.io/docs/model/crtree.html} for an example in Radiant.
\end{Details}
%
\begin{SeeAlso}
\code{\LinkA{crtree}{crtree}} to generate results

\code{\LinkA{summary.crtree}{summary.crtree}} to summarize results

\code{\LinkA{predict.crtree}{predict.crtree}} for prediction
\end{SeeAlso}
%
\begin{Examples}
\begin{ExampleCode}
result <- crtree(titanic, "survived", c("pclass", "sex"), lev = "Yes")
plot(result)
result <- crtree(diamonds, "price", c("carat", "clarity", "cut"))
plot(result, plots = "prune")
result <- crtree(dvd, "buy", c("coupon", "purch", "last"), cp = .01)
plot(result, plots = "imp")

\end{ExampleCode}
\end{Examples}
\HeaderA{plot.dtree}{Plot method for the dtree function}{plot.dtree}
%
\begin{Description}
Plot method for the dtree function
\end{Description}
%
\begin{Usage}
\begin{verbatim}
## S3 method for class 'dtree'
plot(
  x,
  symbol = "$",
  dec = 2,
  final = FALSE,
  orient = "LR",
  width = "900px",
  ...
)
\end{verbatim}
\end{Usage}
%
\begin{Arguments}
\begin{ldescription}
\item[\code{x}] Return value from \code{\LinkA{dtree}{dtree}}

\item[\code{symbol}] Monetary symbol to use (\$ is the default)

\item[\code{dec}] Decimal places to round results to

\item[\code{final}] If TRUE plot the decision tree solution, else the initial decision tree

\item[\code{orient}] Plot orientation: LR for vertical and TD for horizontal

\item[\code{width}] Plot width in pixels (default is "900px")

\item[\code{...}] further arguments passed to or from other methods
\end{ldescription}
\end{Arguments}
%
\begin{Details}
See \url{https://radiant-rstats.github.io/docs/model/dtree.html} for an example in Radiant
\end{Details}
%
\begin{SeeAlso}
\code{\LinkA{dtree}{dtree}} to generate the result

\code{\LinkA{summary.dtree}{summary.dtree}} to summarize results

\code{\LinkA{sensitivity.dtree}{sensitivity.dtree}} to plot results
\end{SeeAlso}
%
\begin{Examples}
\begin{ExampleCode}
dtree(movie_contract, opt = "max") %>% plot()
dtree(movie_contract, opt = "max") %>% plot(final = TRUE, orient = "TD")

\end{ExampleCode}
\end{Examples}
\HeaderA{plot.evalbin}{Plot method for the evalbin function}{plot.evalbin}
%
\begin{Description}
Plot method for the evalbin function
\end{Description}
%
\begin{Usage}
\begin{verbatim}
## S3 method for class 'evalbin'
plot(
  x,
  plots = c("lift", "gains"),
  size = 13,
  shiny = FALSE,
  custom = FALSE,
  ...
)
\end{verbatim}
\end{Usage}
%
\begin{Arguments}
\begin{ldescription}
\item[\code{x}] Return value from \code{\LinkA{evalbin}{evalbin}}

\item[\code{plots}] Plots to return

\item[\code{size}] Font size used

\item[\code{shiny}] Did the function call originate inside a shiny app

\item[\code{custom}] Logical (TRUE, FALSE) to indicate if ggplot object (or list of ggplot objects) should be returned. This option can be used to customize plots (e.g., add a title, change x and y labels, etc.). See examples and \url{https://ggplot2.tidyverse.org} for options.

\item[\code{...}] further arguments passed to or from other methods
\end{ldescription}
\end{Arguments}
%
\begin{Details}
See \url{https://radiant-rstats.github.io/docs/model/evalbin.html} for an example in Radiant
\end{Details}
%
\begin{SeeAlso}
\code{\LinkA{evalbin}{evalbin}} to generate results

\code{\LinkA{summary.evalbin}{summary.evalbin}} to summarize results
\end{SeeAlso}
%
\begin{Examples}
\begin{ExampleCode}
data.frame(buy = dvd$buy, pred1 = runif(20000), pred2 = ifelse(dvd$buy == "yes", 1, 0)) %>%
  evalbin(c("pred1", "pred2"), "buy") %>%
  plot()
\end{ExampleCode}
\end{Examples}
\HeaderA{plot.evalreg}{Plot method for the evalreg function}{plot.evalreg}
%
\begin{Description}
Plot method for the evalreg function
\end{Description}
%
\begin{Usage}
\begin{verbatim}
## S3 method for class 'evalreg'
plot(x, vars = c("Rsq", "RMSE", "MAE"), ...)
\end{verbatim}
\end{Usage}
%
\begin{Arguments}
\begin{ldescription}
\item[\code{x}] Return value from \code{\LinkA{evalreg}{evalreg}}

\item[\code{vars}] Measures to plot, i.e., one or more of "Rsq", "RMSE", "MAE"

\item[\code{...}] further arguments passed to or from other methods
\end{ldescription}
\end{Arguments}
%
\begin{Details}
See \url{https://radiant-rstats.github.io/docs/model/evalreg.html} for an example in Radiant
\end{Details}
%
\begin{SeeAlso}
\code{\LinkA{evalreg}{evalreg}} to generate results

\code{\LinkA{summary.evalreg}{summary.evalreg}} to summarize results
\end{SeeAlso}
%
\begin{Examples}
\begin{ExampleCode}
data.frame(price = diamonds$price, pred1 = rnorm(3000), pred2 = diamonds$price) %>%
  evalreg(pred = c("pred1", "pred2"), "price") %>%
  plot()

\end{ExampleCode}
\end{Examples}
\HeaderA{plot.gbt}{Plot method for the gbt function}{plot.gbt}
%
\begin{Description}
Plot method for the gbt function
\end{Description}
%
\begin{Usage}
\begin{verbatim}
## S3 method for class 'gbt'
plot(
  x,
  plots = "",
  nrobs = Inf,
  incl = NULL,
  incl_int = NULL,
  hline = TRUE,
  pdp_range = c(0.025, 0.975),
  minq = NULL,
  maxq = NULL,
  shiny = FALSE,
  custom = FALSE,
  ...
)
\end{verbatim}
\end{Usage}
%
\begin{Arguments}
\begin{ldescription}
\item[\code{x}] Return value from \code{\LinkA{gbt}{gbt}}

\item[\code{plots}] Plots to produce for the specified Gradient Boosted Tree model. Use "" to avoid showing any plots (default). Options are ...

\item[\code{nrobs}] Number of data points to show in scatter plots (-1 for all)

\item[\code{incl}] Which variables to include in a coefficient plot or PDP plot

\item[\code{incl\_int}] Which interactions to investigate in PDP plots

\item[\code{hline}] Add a dashed horizontal line at the mean response (TRUE/FALSE or a numeric value)

\item[\code{pdp\_range}] Numeric vector \code{c(lo, hi)} giving the percentile range for PDP/Prediction plot x-axes (default \code{c(0.025, 0.975)})

\item[\code{minq}] Deprecated. Use \code{pdp\_range[1]} instead

\item[\code{maxq}] Deprecated. Use \code{pdp\_range[2]} instead

\item[\code{shiny}] Did the function call originate inside a shiny app

\item[\code{custom}] Logical (TRUE, FALSE) to indicate if ggplot object (or list of ggplot objects) should be returned.
This option can be used to customize plots (e.g., add a title, change x and y labels, etc.).
See examples and \url{https://ggplot2.tidyverse.org} for options.

\item[\code{...}] further arguments passed to or from other methods
\end{ldescription}
\end{Arguments}
%
\begin{Details}
See \url{https://radiant-rstats.github.io/docs/model/gbt.html} for an example in Radiant
\end{Details}
%
\begin{SeeAlso}
\code{\LinkA{gbt}{gbt}} to generate results

\code{\LinkA{summary.gbt}{summary.gbt}} to summarize results

\code{\LinkA{predict.gbt}{predict.gbt}} for prediction
\end{SeeAlso}
%
\begin{Examples}
\begin{ExampleCode}
result <- gbt(
  titanic, "survived", c("pclass", "sex"),
  early_stopping_rounds = 0, nthread = 1
)
plot(result)

\end{ExampleCode}
\end{Examples}
\HeaderA{plot.logistic}{Plot method for the logistic function}{plot.logistic}
%
\begin{Description}
Plot method for the logistic function
\end{Description}
%
\begin{Usage}
\begin{verbatim}
## S3 method for class 'logistic'
plot(
  x,
  plots = "coef",
  conf_lev = 0.95,
  intercept = FALSE,
  incl = NULL,
  excl = NULL,
  incl_int = NULL,
  nrobs = -1,
  hline = TRUE,
  pdp_range = c(0.025, 0.975),
  minq = NULL,
  maxq = NULL,
  shiny = FALSE,
  custom = FALSE,
  ...
)
\end{verbatim}
\end{Usage}
%
\begin{Arguments}
\begin{ldescription}
\item[\code{x}] Return value from \code{\LinkA{logistic}{logistic}}

\item[\code{plots}] Plots to produce for the specified GLM model. Use "" to avoid showing any plots (default). "dist" shows histograms (or frequency bar plots) of all variables in the model. "scatter" shows scatter plots (or box plots for factors) for the response variable with each explanatory variable. "coef" provides a coefficient plot and "influence" shows (potentially) influential observations

\item[\code{conf\_lev}] Confidence level to use for coefficient and odds confidence intervals (.95 is the default)

\item[\code{intercept}] Include the intercept in the coefficient plot (TRUE or FALSE). FALSE is the default

\item[\code{incl}] Which variables to include in a coefficient plot

\item[\code{excl}] Which variables to exclude in a coefficient plot

\item[\code{incl\_int}] Which interactions to investigate in PDP plots

\item[\code{nrobs}] Number of data points to show in scatter plots (-1 for all)

\item[\code{hline}] Add a dashed horizontal line at the mean response (TRUE/FALSE or a numeric value)

\item[\code{pdp\_range}] Numeric vector \code{c(lo, hi)} giving the percentile range for PDP/Prediction plot x-axes (default \code{c(0.025, 0.975)})

\item[\code{minq}] Deprecated. Use \code{pdp\_range[1]} instead

\item[\code{maxq}] Deprecated. Use \code{pdp\_range[2]} instead

\item[\code{shiny}] Did the function call originate inside a shiny app

\item[\code{custom}] Logical (TRUE, FALSE) to indicate if ggplot object (or list of ggplot objects) should be returned. This option can be used to customize plots (e.g., add a title, change x and y labels, etc.). See examples and \url{https://ggplot2.tidyverse.org} for options.

\item[\code{...}] further arguments passed to or from other methods
\end{ldescription}
\end{Arguments}
%
\begin{Details}
See \url{https://radiant-rstats.github.io/docs/model/logistic.html} for an example in Radiant
\end{Details}
%
\begin{SeeAlso}
\code{\LinkA{logistic}{logistic}} to generate results

\code{\LinkA{plot.logistic}{plot.logistic}} to plot results

\code{\LinkA{predict.logistic}{predict.logistic}} to generate predictions

\code{\LinkA{plot.model.predict}{plot.model.predict}} to plot prediction output
\end{SeeAlso}
%
\begin{Examples}
\begin{ExampleCode}
result <- logistic(titanic, "survived", c("pclass", "sex"), lev = "Yes")
plot(result, plots = "coef")
\end{ExampleCode}
\end{Examples}
\HeaderA{plot.mnl}{Plot method for the mnl function}{plot.mnl}
%
\begin{Description}
Plot method for the mnl function
\end{Description}
%
\begin{Usage}
\begin{verbatim}
## S3 method for class 'mnl'
plot(
  x,
  plots = "coef",
  conf_lev = 0.95,
  intercept = FALSE,
  nrobs = -1,
  shiny = FALSE,
  custom = FALSE,
  ...
)
\end{verbatim}
\end{Usage}
%
\begin{Arguments}
\begin{ldescription}
\item[\code{x}] Return value from \code{\LinkA{mnl}{mnl}}

\item[\code{plots}] Plots to produce for the specified MNL model. Use "" to avoid showing any plots (default). "dist" shows histograms (or frequency bar plots) of all variables in the model. "scatter" shows scatter plots (or box plots for factors) for the response variable with each explanatory variable. "coef" provides a coefficient plot

\item[\code{conf\_lev}] Confidence level to use for coefficient and relative risk ratios (RRRs) intervals (.95 is the default)

\item[\code{intercept}] Include the intercept in the coefficient plot (TRUE or FALSE). FALSE is the default

\item[\code{nrobs}] Number of data points to show in scatter plots (-1 for all)

\item[\code{shiny}] Did the function call originate inside a shiny app

\item[\code{custom}] Logical (TRUE, FALSE) to indicate if ggplot object (or list of ggplot objects) should be returned. This option can be used to customize plots (e.g., add a title, change x and y labels, etc.). See examples and \url{https://ggplot2.tidyverse.org} for options.

\item[\code{...}] further arguments passed to or from other methods
\end{ldescription}
\end{Arguments}
%
\begin{Details}
See \url{https://radiant-rstats.github.io/docs/model/mnl.html} for an example in Radiant
\end{Details}
%
\begin{SeeAlso}
\code{\LinkA{mnl}{mnl}} to generate results

\code{\LinkA{predict.mnl}{predict.mnl}} to generate predictions

\code{\LinkA{plot.model.predict}{plot.model.predict}} to plot prediction output
\end{SeeAlso}
%
\begin{Examples}
\begin{ExampleCode}
result <- mnl(
  ketchup,
  rvar = "choice",
  evar = c("price.heinz28", "price.heinz32", "price.heinz41", "price.hunts32"),
  lev = "heinz28"
)
plot(result, plots = "coef")

\end{ExampleCode}
\end{Examples}
\HeaderA{plot.mnl.predict}{Plot method for mnl.predict function}{plot.mnl.predict}
%
\begin{Description}
Plot method for mnl.predict function
\end{Description}
%
\begin{Usage}
\begin{verbatim}
## S3 method for class 'mnl.predict'
plot(x, xvar = "", facet_row = ".", facet_col = ".", color = ".class", ...)
\end{verbatim}
\end{Usage}
%
\begin{Arguments}
\begin{ldescription}
\item[\code{x}] Return value from predict function predict.mnl

\item[\code{xvar}] Variable to display along the X-axis of the plot

\item[\code{facet\_row}] Create vertically arranged subplots for each level of the selected factor variable

\item[\code{facet\_col}] Create horizontally arranged subplots for each level of the selected factor variable

\item[\code{color}] Adds color to a scatter plot to generate a heat map. For a line plot one line is created for each group and each is assigned a different color

\item[\code{...}] further arguments passed to or from other methods
\end{ldescription}
\end{Arguments}
%
\begin{SeeAlso}
\code{\LinkA{predict.mnl}{predict.mnl}} to generate predictions
\end{SeeAlso}
%
\begin{Examples}
\begin{ExampleCode}
result <- mnl(
  ketchup,
  rvar = "choice",
  evar = c("price.heinz28", "price.heinz32", "price.heinz41", "price.hunts32"),
  lev = "heinz28"
)
pred <- predict(result, pred_cmd = "price.heinz28 = seq(3, 5, 0.1)")
plot(pred, xvar = "price.heinz28")

\end{ExampleCode}
\end{Examples}
\HeaderA{plot.model.predict}{Plot method for model.predict functions}{plot.model.predict}
%
\begin{Description}
Plot method for model.predict functions
\end{Description}
%
\begin{Usage}
\begin{verbatim}
## S3 method for class 'model.predict'
plot(
  x,
  xvar = "",
  facet_row = ".",
  facet_col = ".",
  color = "none",
  conf_lev = 0.95,
  ...
)
\end{verbatim}
\end{Usage}
%
\begin{Arguments}
\begin{ldescription}
\item[\code{x}] Return value from predict functions (e.g., predict.regress)

\item[\code{xvar}] Variable to display along the X-axis of the plot

\item[\code{facet\_row}] Create vertically arranged subplots for each level of the selected factor variable

\item[\code{facet\_col}] Create horizontally arranged subplots for each level of the selected factor variable

\item[\code{color}] Adds color to a scatter plot to generate a heat map. For a line plot one line is created for each group and each is assigned a different color

\item[\code{conf\_lev}] Confidence level to use for prediction intervals (.95 is the default)

\item[\code{...}] further arguments passed to or from other methods
\end{ldescription}
\end{Arguments}
%
\begin{SeeAlso}
\code{\LinkA{predict.regress}{predict.regress}} to generate predictions

\code{\LinkA{predict.logistic}{predict.logistic}} to generate predictions
\end{SeeAlso}
%
\begin{Examples}
\begin{ExampleCode}
regress(diamonds, "price", c("carat", "clarity")) %>%
  predict(pred_cmd = "carat = 1:10") %>%
  plot(xvar = "carat")
logistic(titanic, "survived", c("pclass", "sex", "age"), lev = "Yes") %>%
  predict(pred_cmd = c("pclass = levels(pclass)", "sex = levels(sex)", "age = 0:100")) %>%
  plot(xvar = "age", color = "sex", facet_col = "pclass")

\end{ExampleCode}
\end{Examples}
\HeaderA{plot.nb}{Plot method for the nb function}{plot.nb}
%
\begin{Description}
Plot method for the nb function
\end{Description}
%
\begin{Usage}
\begin{verbatim}
## S3 method for class 'nb'
plot(x, plots = "correlations", lev = "All levels", nrobs = 1000, ...)
\end{verbatim}
\end{Usage}
%
\begin{Arguments}
\begin{ldescription}
\item[\code{x}] Return value from \code{\LinkA{nb}{nb}}

\item[\code{plots}] Plots to produce for the specified model. Use "" to avoid showing any plots. Use "vimp" for variable importance or "correlations" to examine conditional independence

\item[\code{lev}] The level(s) in the response variable used as the basis for plots (defaults to "All levels")

\item[\code{nrobs}] Number of data points to show in scatter plots (-1 for all)

\item[\code{...}] further arguments passed to or from other methods
\end{ldescription}
\end{Arguments}
%
\begin{Details}
See \url{https://radiant-rstats.github.io/docs/model/nb.html} for an example in Radiant
\end{Details}
%
\begin{SeeAlso}
\code{\LinkA{nb}{nb}} to generate results

\code{\LinkA{summary.nb}{summary.nb}} to summarize results

\code{\LinkA{predict.nb}{predict.nb}} for prediction
\end{SeeAlso}
%
\begin{Examples}
\begin{ExampleCode}
result <- nb(titanic, "survived", c("pclass", "sex"))
plot(result)
result <- nb(titanic, "pclass", c("sex", "age"))
plot(result)

\end{ExampleCode}
\end{Examples}
\HeaderA{plot.nb.predict}{Plot method for nb.predict function}{plot.nb.predict}
%
\begin{Description}
Plot method for nb.predict function
\end{Description}
%
\begin{Usage}
\begin{verbatim}
## S3 method for class 'nb.predict'
plot(x, xvar = "", facet_row = ".", facet_col = ".", color = ".class", ...)
\end{verbatim}
\end{Usage}
%
\begin{Arguments}
\begin{ldescription}
\item[\code{x}] Return value from predict function predict.nb

\item[\code{xvar}] Variable to display along the X-axis of the plot

\item[\code{facet\_row}] Create vertically arranged subplots for each level of the selected factor variable

\item[\code{facet\_col}] Create horizontally arranged subplots for each level of the selected factor variable

\item[\code{color}] Adds color to a scatter plot to generate a heat map. For a line plot one line is created for each group and each is assigned a different color

\item[\code{...}] further arguments passed to or from other methods
\end{ldescription}
\end{Arguments}
%
\begin{SeeAlso}
\code{\LinkA{predict.nb}{predict.nb}} to generate predictions
\end{SeeAlso}
%
\begin{Examples}
\begin{ExampleCode}
result <- nb(titanic, "survived", c("pclass", "sex", "age"))
pred <- predict(
  result,
  pred_cmd = c("pclass = levels(pclass)", "sex = levels(sex)", "age = seq(0, 100, 20)")
)
plot(pred, xvar = "age", facet_col = "sex", facet_row = "pclass")
pred <- predict(result, pred_data = titanic)
plot(pred, xvar = "age", facet_col = "sex")

\end{ExampleCode}
\end{Examples}
\HeaderA{plot.nn}{Plot method for the nn function}{plot.nn}
%
\begin{Description}
Plot method for the nn function
\end{Description}
%
\begin{Usage}
\begin{verbatim}
## S3 method for class 'nn'
plot(
  x,
  plots = "vip",
  size = 12,
  pad_x = 0.9,
  nrobs = -1,
  incl = NULL,
  incl_int = NULL,
  hline = TRUE,
  pdp_range = c(0.025, 0.975),
  minq = NULL,
  maxq = NULL,
  shiny = FALSE,
  custom = FALSE,
  ...
)
\end{verbatim}
\end{Usage}
%
\begin{Arguments}
\begin{ldescription}
\item[\code{x}] Return value from \code{\LinkA{nn}{nn}}

\item[\code{plots}] Plots to produce for the specified Neural Network model. Use "" to avoid showing any plots (default). Options are "olden" or "garson" for importance plots, or "net" to depict the network structure

\item[\code{size}] Font size used

\item[\code{pad\_x}] Padding for explanatory variable labels in the network plot. Default value is 0.9, smaller numbers (e.g., 0.5) increase the amount of padding

\item[\code{nrobs}] Number of data points to show in dashboard scatter plots (-1 for all)

\item[\code{incl}] Which variables to include in a coefficient plot or PDP plot

\item[\code{incl\_int}] Which interactions to investigate in PDP plots

\item[\code{hline}] Add a dashed horizontal line at the mean response (TRUE/FALSE or a numeric value)

\item[\code{pdp\_range}] Numeric vector \code{c(lo, hi)} giving the percentile range for PDP/Prediction plot x-axes (default \code{c(0.025, 0.975)})

\item[\code{minq}] Deprecated. Use \code{pdp\_range[1]} instead

\item[\code{maxq}] Deprecated. Use \code{pdp\_range[2]} instead

\item[\code{shiny}] Did the function call originate inside a shiny app

\item[\code{custom}] Logical (TRUE, FALSE) to indicate if ggplot object (or list of ggplot objects) should be returned. This option can be used to customize plots (e.g., add a title, change x and y labels, etc.). See examples and \url{https://ggplot2.tidyverse.org} for options.

\item[\code{...}] further arguments passed to or from other methods
\end{ldescription}
\end{Arguments}
%
\begin{Details}
See \url{https://radiant-rstats.github.io/docs/model/nn.html} for an example in Radiant
\end{Details}
%
\begin{SeeAlso}
\code{\LinkA{nn}{nn}} to generate results

\code{\LinkA{summary.nn}{summary.nn}} to summarize results

\code{\LinkA{predict.nn}{predict.nn}} for prediction
\end{SeeAlso}
%
\begin{Examples}
\begin{ExampleCode}
result <- nn(titanic, "survived", c("pclass", "sex"), lev = "Yes")
plot(result, plots = "net")
plot(result, plots = "olden")
\end{ExampleCode}
\end{Examples}
\HeaderA{plot.regress}{Plot method for the regress function}{plot.regress}
%
\begin{Description}
Plot method for the regress function
\end{Description}
%
\begin{Usage}
\begin{verbatim}
## S3 method for class 'regress'
plot(
  x,
  plots = "",
  lines = "",
  conf_lev = 0.95,
  intercept = FALSE,
  incl = NULL,
  excl = NULL,
  incl_int = NULL,
  nrobs = -1,
  hline = TRUE,
  pdp_range = c(0.025, 0.975),
  minq = NULL,
  maxq = NULL,
  shiny = FALSE,
  custom = FALSE,
  ...
)
\end{verbatim}
\end{Usage}
%
\begin{Arguments}
\begin{ldescription}
\item[\code{x}] Return value from \code{\LinkA{regress}{regress}}

\item[\code{plots}] Regression plots to produce for the specified regression model. Enter "" to avoid showing any plots (default). "dist" to shows histograms (or frequency bar plots) of all variables in the model. "correlations" for a visual representation of the correlation matrix selected variables. "scatter" to show scatter plots (or box plots for factors) for the response variable with each explanatory variable. "dashboard" for a series of six plots that can be used to evaluate model fit visually. "resid\_pred" to plot the explanatory variables against the model residuals. "coef" for a coefficient plot with adjustable confidence intervals and "influence" to show (potentially) influential observations

\item[\code{lines}] Optional lines to include in the select plot. "line" to include a line through a scatter plot. "loess" to include a polynomial regression fit line. To include both use c("line", "loess")

\item[\code{conf\_lev}] Confidence level used to estimate confidence intervals (.95 is the default)

\item[\code{intercept}] Include the intercept in the coefficient plot (TRUE, FALSE). FALSE is the default

\item[\code{incl}] Which variables to include in a coefficient plot or PDP plot

\item[\code{excl}] Which variables to exclude in a coefficient plot

\item[\code{incl\_int}] Which interactions to investigate in PDP plots

\item[\code{nrobs}] Number of data points to show in scatter plots (-1 for all)

\item[\code{hline}] Add a dashed horizontal line at the mean response (TRUE/FALSE or a numeric value)

\item[\code{pdp\_range}] Numeric vector \code{c(lo, hi)} giving the percentile range for PDP/Prediction plot x-axes (default \code{c(0.025, 0.975)})

\item[\code{minq}] Deprecated. Use \code{pdp\_range[1]} instead

\item[\code{maxq}] Deprecated. Use \code{pdp\_range[2]} instead

\item[\code{shiny}] Did the function call originate inside a shiny app

\item[\code{custom}] Logical (TRUE, FALSE) to indicate if ggplot object (or list of ggplot objects) should be returned. This option can be used to customize plots (e.g., add a title, change x and y labels, etc.). See examples and \url{https://ggplot2.tidyverse.org} for options.

\item[\code{...}] further arguments passed to or from other methods
\end{ldescription}
\end{Arguments}
%
\begin{Details}
See \url{https://radiant-rstats.github.io/docs/model/regress.html} for an example in Radiant
\end{Details}
%
\begin{SeeAlso}
\code{\LinkA{regress}{regress}} to generate the results

\code{\LinkA{summary.regress}{summary.regress}} to summarize results

\code{\LinkA{predict.regress}{predict.regress}} to generate predictions
\end{SeeAlso}
%
\begin{Examples}
\begin{ExampleCode}
result <- regress(diamonds, "price", c("carat", "clarity"))
plot(result, plots = "coef", conf_lev = .99, intercept = TRUE)
## Not run: 
plot(result, plots = "dist")
plot(result, plots = "scatter", lines = c("line", "loess"))
plot(result, plots = "resid_pred", lines = "line")
plot(result, plots = "dashboard", lines = c("line", "loess"))

## End(Not run)
\end{ExampleCode}
\end{Examples}
\HeaderA{plot.repeater}{Plot repeated simulation}{plot.repeater}
%
\begin{Description}
Plot repeated simulation
\end{Description}
%
\begin{Usage}
\begin{verbatim}
## S3 method for class 'repeater'
plot(x, bins = 20, shiny = FALSE, custom = FALSE, ...)
\end{verbatim}
\end{Usage}
%
\begin{Arguments}
\begin{ldescription}
\item[\code{x}] Return value from \code{\LinkA{repeater}{repeater}}

\item[\code{bins}] Number of bins used for histograms (1 - 50)

\item[\code{shiny}] Did the function call originate inside a shiny app

\item[\code{custom}] Logical (TRUE, FALSE) to indicate if ggplot object (or list of ggplot objects) should be returned. This option can be used to customize plots (e.g., add a title, change x and y labels, etc.). See examples and \url{https://ggplot2.tidyverse.org} for options.

\item[\code{...}] further arguments passed to or from other methods
\end{ldescription}
\end{Arguments}
%
\begin{SeeAlso}
\code{\LinkA{repeater}{repeater}} to run a repeated simulation

\code{\LinkA{summary.repeater}{summary.repeater}} to summarize results from repeated simulation
\end{SeeAlso}
\HeaderA{plot.rforest}{Plot method for the rforest function}{plot.rforest}
%
\begin{Description}
Plot method for the rforest function
\end{Description}
%
\begin{Usage}
\begin{verbatim}
## S3 method for class 'rforest'
plot(
  x,
  plots = "",
  nrobs = Inf,
  incl = NULL,
  incl_int = NULL,
  hline = TRUE,
  pdp_range = c(0.025, 0.975),
  minq = NULL,
  maxq = NULL,
  shiny = FALSE,
  custom = FALSE,
  ...
)
\end{verbatim}
\end{Usage}
%
\begin{Arguments}
\begin{ldescription}
\item[\code{x}] Return value from \code{\LinkA{rforest}{rforest}}

\item[\code{plots}] Plots to produce for the specified Random Forest model. Use "" to avoid showing any plots (default). Options are ...

\item[\code{nrobs}] Number of data points to show in dashboard scatter plots (-1 for all)

\item[\code{incl}] Which variables to include in PDP or Prediction plots

\item[\code{incl\_int}] Which interactions to investigate in PDP or Prediction plots

\item[\code{hline}] Add a dashed horizontal line at the mean response (TRUE/FALSE or a numeric value)

\item[\code{pdp\_range}] Numeric vector \code{c(lo, hi)} giving the percentile range for PDP/Prediction plot x-axes (default \code{c(0.025, 0.975)})

\item[\code{minq}] Deprecated. Use \code{pdp\_range[1]} instead

\item[\code{maxq}] Deprecated. Use \code{pdp\_range[2]} instead

\item[\code{shiny}] Did the function call originate inside a shiny app

\item[\code{custom}] Logical (TRUE, FALSE) to indicate if ggplot object (or list of ggplot objects) should be returned.
This option can be used to customize plots (e.g., add a title, change x and y labels, etc.). See examples
and \url{https://ggplot2.tidyverse.org} for options.

\item[\code{...}] further arguments passed to or from other methods
\end{ldescription}
\end{Arguments}
%
\begin{Details}
See \url{https://radiant-rstats.github.io/docs/model/rforest.html} for an example in Radiant
\end{Details}
%
\begin{SeeAlso}
\code{\LinkA{rforest}{rforest}} to generate results

\code{\LinkA{summary.rforest}{summary.rforest}} to summarize results

\code{\LinkA{predict.rforest}{predict.rforest}} for prediction
\end{SeeAlso}
%
\begin{Examples}
\begin{ExampleCode}
result <- rforest(titanic, "survived", c("pclass", "sex"), lev = "Yes")

\end{ExampleCode}
\end{Examples}
\HeaderA{plot.rforest.predict}{Plot method for rforest.predict function}{plot.rforest.predict}
%
\begin{Description}
Plot method for rforest.predict function
\end{Description}
%
\begin{Usage}
\begin{verbatim}
## S3 method for class 'rforest.predict'
plot(x, xvar = "", facet_row = ".", facet_col = ".", color = "none", ...)
\end{verbatim}
\end{Usage}
%
\begin{Arguments}
\begin{ldescription}
\item[\code{x}] Return value from predict function predict.rforest

\item[\code{xvar}] Variable to display along the X-axis of the plot

\item[\code{facet\_row}] Create vertically arranged subplots for each level of the selected factor variable

\item[\code{facet\_col}] Create horizontally arranged subplots for each level of the selected factor variable

\item[\code{color}] Adds color to a scatter plot to generate a heat map. For a line plot one line is created for each group and each is assigned a different color

\item[\code{...}] further arguments passed to or from other methods
\end{ldescription}
\end{Arguments}
%
\begin{SeeAlso}
\code{\LinkA{predict.mnl}{predict.mnl}} to generate predictions
\end{SeeAlso}
%
\begin{Examples}
\begin{ExampleCode}
result <- mnl(
  ketchup,
  rvar = "choice",
  evar = c("price.heinz28", "price.heinz32", "price.heinz41", "price.hunts32"),
  lev = "heinz28"
)
pred <- predict(result, pred_cmd = "price.heinz28 = seq(3, 5, 0.1)")
plot(pred, xvar = "price.heinz28")

\end{ExampleCode}
\end{Examples}
\HeaderA{plot.simulater}{Plot method for the simulater function}{plot.simulater}
%
\begin{Description}
Plot method for the simulater function
\end{Description}
%
\begin{Usage}
\begin{verbatim}
## S3 method for class 'simulater'
plot(x, bins = 20, shiny = FALSE, custom = FALSE, ...)
\end{verbatim}
\end{Usage}
%
\begin{Arguments}
\begin{ldescription}
\item[\code{x}] Return value from \code{\LinkA{simulater}{simulater}}

\item[\code{bins}] Number of bins used for histograms (1 - 50)

\item[\code{shiny}] Did the function call originate inside a shiny app

\item[\code{custom}] Logical (TRUE, FALSE) to indicate if ggplot object (or list of ggplot objects) should be returned. This option can be used to customize plots (e.g., add a title, change x and y labels, etc.). See examples and \url{https://ggplot2.tidyverse.org} for options.

\item[\code{...}] further arguments passed to or from other methods
\end{ldescription}
\end{Arguments}
%
\begin{Details}
See \url{https://radiant-rstats.github.io/docs/model/simulater} for an example in Radiant
\end{Details}
%
\begin{SeeAlso}
\code{\LinkA{simulater}{simulater}} to generate the result

\code{\LinkA{summary.simulater}{summary.simulater}} to summarize results
\end{SeeAlso}
%
\begin{Examples}
\begin{ExampleCode}
simdat <- simulater(
  const = "cost 3",
  norm = "demand 2000 1000",
  discrete = "price 5 8 .3 .7",
  form = "profit = demand * (price - cost)",
  seed = 1234
)
plot(simdat, bins = 25)

\end{ExampleCode}
\end{Examples}
\HeaderA{plot.uplift}{Plot method for the uplift function}{plot.uplift}
%
\begin{Description}
Plot method for the uplift function
\end{Description}
%
\begin{Usage}
\begin{verbatim}
## S3 method for class 'uplift'
plot(
  x,
  plots = c("inc_uplift", "uplift"),
  size = 13,
  shiny = FALSE,
  custom = FALSE,
  ...
)
\end{verbatim}
\end{Usage}
%
\begin{Arguments}
\begin{ldescription}
\item[\code{x}] Return value from \code{\LinkA{evalbin}{evalbin}}

\item[\code{plots}] Plots to return

\item[\code{size}] Font size used

\item[\code{shiny}] Did the function call originate inside a shiny app

\item[\code{custom}] Logical (TRUE, FALSE) to indicate if ggplot object (or list of ggplot objects) should be returned. This option can be used to customize plots (e.g., add a title, change x and y labels, etc.). See examples and \url{https://ggplot2.tidyverse.org} for options.

\item[\code{...}] further arguments passed to or from other methods
\end{ldescription}
\end{Arguments}
%
\begin{Details}
See \url{https://radiant-rstats.github.io/docs/model/evalbin.html} for an example in Radiant
\end{Details}
%
\begin{SeeAlso}
\code{\LinkA{evalbin}{evalbin}} to generate results

\code{\LinkA{summary.evalbin}{summary.evalbin}} to summarize results
\end{SeeAlso}
%
\begin{Examples}
\begin{ExampleCode}
data.frame(buy = dvd$buy, pred1 = runif(20000), pred2 = ifelse(dvd$buy == "yes", 1, 0)) %>%
  evalbin(c("pred1", "pred2"), "buy") %>%
  plot()
\end{ExampleCode}
\end{Examples}
\HeaderA{predict.crtree}{Predict method for the crtree function}{predict.crtree}
%
\begin{Description}
Predict method for the crtree function
\end{Description}
%
\begin{Usage}
\begin{verbatim}
## S3 method for class 'crtree'
predict(
  object,
  pred_data = NULL,
  pred_cmd = "",
  conf_lev = 0.95,
  se = FALSE,
  dec = 3,
  envir = parent.frame(),
  ...
)
\end{verbatim}
\end{Usage}
%
\begin{Arguments}
\begin{ldescription}
\item[\code{object}] Return value from \code{\LinkA{crtree}{crtree}}

\item[\code{pred\_data}] Provide the dataframe to generate predictions (e.g., titanic). The dataset must contain all columns used in the estimation

\item[\code{pred\_cmd}] Generate predictions using a command. For example, `pclass = levels(pclass)` would produce predictions for the different levels of factor `pclass`. To add another variable, create a vector of prediction strings, (e.g., c('pclass = levels(pclass)', 'age = seq(0,100,20)')

\item[\code{conf\_lev}] Confidence level used to estimate confidence intervals (.95 is the default)

\item[\code{se}] Logical that indicates if prediction standard errors should be calculated (default = FALSE)

\item[\code{dec}] Number of decimals to show

\item[\code{envir}] Environment to extract data from

\item[\code{...}] further arguments passed to or from other methods
\end{ldescription}
\end{Arguments}
%
\begin{Details}
See \url{https://radiant-rstats.github.io/docs/model/crtree.html} for an example in Radiant
\end{Details}
%
\begin{SeeAlso}
\code{\LinkA{crtree}{crtree}} to generate the result

\code{\LinkA{summary.crtree}{summary.crtree}} to summarize results
\end{SeeAlso}
%
\begin{Examples}
\begin{ExampleCode}
result <- crtree(titanic, "survived", c("pclass", "sex"), lev = "Yes")
predict(result, pred_cmd = "pclass = levels(pclass)")
result <- crtree(titanic, "survived", "pclass", lev = "Yes")
predict(result, pred_data = titanic) %>% head()
\end{ExampleCode}
\end{Examples}
\HeaderA{predict.gbt}{Predict method for the gbt function}{predict.gbt}
%
\begin{Description}
Predict method for the gbt function
\end{Description}
%
\begin{Usage}
\begin{verbatim}
## S3 method for class 'gbt'
predict(
  object,
  pred_data = NULL,
  pred_cmd = "",
  dec = 3,
  envir = parent.frame(),
  ...
)
\end{verbatim}
\end{Usage}
%
\begin{Arguments}
\begin{ldescription}
\item[\code{object}] Return value from \code{\LinkA{gbt}{gbt}}

\item[\code{pred\_data}] Provide the dataframe to generate predictions (e.g., diamonds). The dataset must contain all columns used in the estimation

\item[\code{pred\_cmd}] Generate predictions using a command. For example, `pclass = levels(pclass)` would produce predictions for the different levels of factor `pclass`. To add another variable, create a vector of prediction strings, (e.g., c('pclass = levels(pclass)', 'age = seq(0,100,20)')

\item[\code{dec}] Number of decimals to show

\item[\code{envir}] Environment to extract data from

\item[\code{...}] further arguments passed to or from other methods
\end{ldescription}
\end{Arguments}
%
\begin{Details}
See \url{https://radiant-rstats.github.io/docs/model/gbt.html} for an example in Radiant
\end{Details}
%
\begin{SeeAlso}
\code{\LinkA{gbt}{gbt}} to generate the result

\code{\LinkA{summary.gbt}{summary.gbt}} to summarize results
\end{SeeAlso}
%
\begin{Examples}
\begin{ExampleCode}
result <- gbt(
  titanic, "survived", c("pclass", "sex"),
  early_stopping_rounds = 2, nthread = 1
)
predict(result, pred_cmd = "pclass = levels(pclass)")
result <- gbt(diamonds, "price", "carat:color", type = "regression", nthread = 1)
predict(result, pred_cmd = "carat = 1:3")
predict(result, pred_data = diamonds) %>% head()
\end{ExampleCode}
\end{Examples}
\HeaderA{predict.logistic}{Predict method for the logistic function}{predict.logistic}
%
\begin{Description}
Predict method for the logistic function
\end{Description}
%
\begin{Usage}
\begin{verbatim}
## S3 method for class 'logistic'
predict(
  object,
  pred_data = NULL,
  pred_cmd = "",
  conf_lev = 0.95,
  se = TRUE,
  interval = "confidence",
  dec = 3,
  envir = parent.frame(),
  ...
)
\end{verbatim}
\end{Usage}
%
\begin{Arguments}
\begin{ldescription}
\item[\code{object}] Return value from \code{\LinkA{logistic}{logistic}}

\item[\code{pred\_data}] Provide the dataframe to generate predictions (e.g., titanic). The dataset must contain all columns used in the estimation

\item[\code{pred\_cmd}] Generate predictions using a command. For example, `pclass = levels(pclass)` would produce predictions for the different levels of factor `pclass`. To add another variable, create a vector of prediction strings, (e.g., c('pclass = levels(pclass)', 'age = seq(0,100,20)')

\item[\code{conf\_lev}] Confidence level used to estimate confidence intervals (.95 is the default)

\item[\code{se}] Logical that indicates if prediction standard errors should be calculated (default = FALSE)

\item[\code{interval}] Type of interval calculation ("confidence" or "none"). Set to "none" if se is FALSE

\item[\code{dec}] Number of decimals to show

\item[\code{envir}] Environment to extract data from

\item[\code{...}] further arguments passed to or from other methods
\end{ldescription}
\end{Arguments}
%
\begin{Details}
See \url{https://radiant-rstats.github.io/docs/model/logistic.html} for an example in Radiant
\end{Details}
%
\begin{SeeAlso}
\code{\LinkA{logistic}{logistic}} to generate the result

\code{\LinkA{summary.logistic}{summary.logistic}} to summarize results

\code{\LinkA{plot.logistic}{plot.logistic}} to plot results

\code{\LinkA{plot.model.predict}{plot.model.predict}} to plot prediction output
\end{SeeAlso}
%
\begin{Examples}
\begin{ExampleCode}
result <- logistic(titanic, "survived", c("pclass", "sex"), lev = "Yes")
predict(result, pred_cmd = "pclass = levels(pclass)")
logistic(titanic, "survived", c("pclass", "sex"), lev = "Yes") %>%
  predict(pred_cmd = "sex = c('male','female')")
logistic(titanic, "survived", c("pclass", "sex"), lev = "Yes") %>%
  predict(pred_data = titanic)
\end{ExampleCode}
\end{Examples}
\HeaderA{predict.mnl}{Predict method for the mnl function}{predict.mnl}
%
\begin{Description}
Predict method for the mnl function
\end{Description}
%
\begin{Usage}
\begin{verbatim}
## S3 method for class 'mnl'
predict(
  object,
  pred_data = NULL,
  pred_cmd = "",
  pred_names = "",
  dec = 3,
  envir = parent.frame(),
  ...
)
\end{verbatim}
\end{Usage}
%
\begin{Arguments}
\begin{ldescription}
\item[\code{object}] Return value from \code{\LinkA{mnl}{mnl}}

\item[\code{pred\_data}] Provide the dataframe to generate predictions (e.g., ketchup). The dataset must contain all columns used in the estimation

\item[\code{pred\_cmd}] Generate predictions using a command. For example, `pclass = levels(pclass)` would produce predictions for the different levels of factor `pclass`. To add another variable, create a vector of prediction strings, (e.g., c('pclass = levels(pclass)', 'age = seq(0,100,20)')

\item[\code{pred\_names}] Names for the predictions to be stored. If one name is provided, only the first column of predictions is stored. If empty, the levels in the response variable of the mnl model will be used

\item[\code{dec}] Number of decimals to show

\item[\code{envir}] Environment to extract data from

\item[\code{...}] further arguments passed to or from other methods
\end{ldescription}
\end{Arguments}
%
\begin{Details}
See \url{https://radiant-rstats.github.io/docs/model/mnl.html} for an example in Radiant
\end{Details}
%
\begin{SeeAlso}
\code{\LinkA{mnl}{mnl}} to generate the result

\code{\LinkA{summary.mnl}{summary.mnl}} to summarize results
\end{SeeAlso}
%
\begin{Examples}
\begin{ExampleCode}
result <- mnl(
  ketchup,
  rvar = "choice",
  evar = c("price.heinz28", "price.heinz32", "price.heinz41", "price.hunts32"),
  lev = "heinz28"
)
predict(result, pred_cmd = "price.heinz28 = seq(3, 5, 0.1)")
predict(result, pred_data = slice(ketchup, 1:20))

\end{ExampleCode}
\end{Examples}
\HeaderA{predict.nb}{Predict method for the nb function}{predict.nb}
%
\begin{Description}
Predict method for the nb function
\end{Description}
%
\begin{Usage}
\begin{verbatim}
## S3 method for class 'nb'
predict(
  object,
  pred_data = NULL,
  pred_cmd = "",
  pred_names = "",
  dec = 3,
  envir = parent.frame(),
  ...
)
\end{verbatim}
\end{Usage}
%
\begin{Arguments}
\begin{ldescription}
\item[\code{object}] Return value from \code{\LinkA{nb}{nb}}

\item[\code{pred\_data}] Provide the dataframe to generate predictions (e.g., titanic). The dataset must contain all columns used in the estimation

\item[\code{pred\_cmd}] Generate predictions using a command. For example, `pclass = levels(pclass)` would produce predictions for the different levels of factor `pclass`. To add another variable, create a vector of prediction strings, (e.g., c('pclass = levels(pclass)', 'age = seq(0,100,20)')

\item[\code{pred\_names}] Names for the predictions to be stored. If one name is provided, only the first column of predictions is stored. If empty, the level in the response variable of the nb model will be used

\item[\code{dec}] Number of decimals to show

\item[\code{envir}] Environment to extract data from

\item[\code{...}] further arguments passed to or from other methods
\end{ldescription}
\end{Arguments}
%
\begin{Details}
See \url{https://radiant-rstats.github.io/docs/model/nb.html} for an example in Radiant
\end{Details}
%
\begin{SeeAlso}
\code{\LinkA{nb}{nb}} to generate the result

\code{\LinkA{summary.nb}{summary.nb}} to summarize results
\end{SeeAlso}
%
\begin{Examples}
\begin{ExampleCode}
result <- nb(titanic, "survived", c("pclass", "sex", "age"))
predict(result, pred_data = titanic)
predict(result, pred_data = titanic, pred_names = c("Yes", "No"))
predict(result, pred_cmd = "pclass = levels(pclass)")
result <- nb(titanic, "pclass", c("survived", "sex", "age"))
predict(result, pred_data = titanic)
predict(result, pred_data = titanic, pred_names = c("1st", "2nd", "3rd"))
predict(result, pred_data = titanic, pred_names = "")

\end{ExampleCode}
\end{Examples}
\HeaderA{predict.nn}{Predict method for the nn function}{predict.nn}
%
\begin{Description}
Predict method for the nn function
\end{Description}
%
\begin{Usage}
\begin{verbatim}
## S3 method for class 'nn'
predict(
  object,
  pred_data = NULL,
  pred_cmd = "",
  dec = 3,
  envir = parent.frame(),
  ...
)
\end{verbatim}
\end{Usage}
%
\begin{Arguments}
\begin{ldescription}
\item[\code{object}] Return value from \code{\LinkA{nn}{nn}}

\item[\code{pred\_data}] Provide the dataframe to generate predictions (e.g., diamonds). The dataset must contain all columns used in the estimation

\item[\code{pred\_cmd}] Generate predictions using a command. For example, `pclass = levels(pclass)` would produce predictions for the different levels of factor `pclass`. To add another variable, create a vector of prediction strings, (e.g., c('pclass = levels(pclass)', 'age = seq(0,100,20)')

\item[\code{dec}] Number of decimals to show

\item[\code{envir}] Environment to extract data from

\item[\code{...}] further arguments passed to or from other methods
\end{ldescription}
\end{Arguments}
%
\begin{Details}
See \url{https://radiant-rstats.github.io/docs/model/nn.html} for an example in Radiant
\end{Details}
%
\begin{SeeAlso}
\code{\LinkA{nn}{nn}} to generate the result

\code{\LinkA{summary.nn}{summary.nn}} to summarize results
\end{SeeAlso}
%
\begin{Examples}
\begin{ExampleCode}
result <- nn(titanic, "survived", c("pclass", "sex"), lev = "Yes")
predict(result, pred_cmd = "pclass = levels(pclass)")
result <- nn(diamonds, "price", "carat:color", type = "regression")
predict(result, pred_cmd = "carat = 1:3")
predict(result, pred_data = diamonds) %>% head()
\end{ExampleCode}
\end{Examples}
\HeaderA{predict.regress}{Predict method for the regress function}{predict.regress}
%
\begin{Description}
Predict method for the regress function
\end{Description}
%
\begin{Usage}
\begin{verbatim}
## S3 method for class 'regress'
predict(
  object,
  pred_data = NULL,
  pred_cmd = "",
  conf_lev = 0.95,
  se = TRUE,
  interval = "confidence",
  dec = 3,
  envir = parent.frame(),
  ...
)
\end{verbatim}
\end{Usage}
%
\begin{Arguments}
\begin{ldescription}
\item[\code{object}] Return value from \code{\LinkA{regress}{regress}}

\item[\code{pred\_data}] Provide the dataframe to generate predictions (e.g., diamonds). The dataset must contain all columns used in the estimation

\item[\code{pred\_cmd}] Command used to generate data for prediction

\item[\code{conf\_lev}] Confidence level used to estimate confidence intervals (.95 is the default)

\item[\code{se}] Logical that indicates if prediction standard errors should be calculated (default = FALSE)

\item[\code{interval}] Type of interval calculation ("confidence" or "prediction"). Set to "none" if se is FALSE

\item[\code{dec}] Number of decimals to show

\item[\code{envir}] Environment to extract data from

\item[\code{...}] further arguments passed to or from other methods
\end{ldescription}
\end{Arguments}
%
\begin{Details}
See \url{https://radiant-rstats.github.io/docs/model/regress.html} for an example in Radiant
\end{Details}
%
\begin{SeeAlso}
\code{\LinkA{regress}{regress}} to generate the result

\code{\LinkA{summary.regress}{summary.regress}} to summarize results

\code{\LinkA{plot.regress}{plot.regress}} to plot results
\end{SeeAlso}
%
\begin{Examples}
\begin{ExampleCode}
result <- regress(diamonds, "price", c("carat", "clarity"))
predict(result, pred_cmd = "carat = 1:10")
predict(result, pred_cmd = "clarity = levels(clarity)")
result <- regress(diamonds, "price", c("carat", "clarity"), int = "carat:clarity")
predict(result, pred_data = diamonds) %>% head()

\end{ExampleCode}
\end{Examples}
\HeaderA{predict.rforest}{Predict method for the rforest function}{predict.rforest}
%
\begin{Description}
Predict method for the rforest function
\end{Description}
%
\begin{Usage}
\begin{verbatim}
## S3 method for class 'rforest'
predict(
  object,
  pred_data = NULL,
  pred_cmd = "",
  pred_names = "",
  OOB = NULL,
  dec = 3,
  envir = parent.frame(),
  ...
)
\end{verbatim}
\end{Usage}
%
\begin{Arguments}
\begin{ldescription}
\item[\code{object}] Return value from \code{\LinkA{rforest}{rforest}}

\item[\code{pred\_data}] Provide the dataframe to generate predictions (e.g., diamonds). The dataset must contain all columns used in the estimation

\item[\code{pred\_cmd}] Generate predictions using a command. For example, `pclass = levels(pclass)` would produce predictions for the different
levels of factor `pclass`. To add another variable, create a vector of prediction strings, (e.g., c('pclass = levels(pclass)', 'age = seq(0,100,20)')

\item[\code{pred\_names}] Names for the predictions to be stored. If one name is provided, only the first column of predictions is stored. If empty, the levels
in the response variable of the rforest model will be used

\item[\code{OOB}] Use Out-Of-Bag predictions (TRUE or FALSE). Relevant when evaluating predictions for the training sample. If set to NULL, datasets will be compared
to determine if OOB predictions should be used

\item[\code{dec}] Number of decimals to show

\item[\code{envir}] Environment to extract data from

\item[\code{...}] further arguments passed to or from other methods
\end{ldescription}
\end{Arguments}
%
\begin{Details}
See \url{https://radiant-rstats.github.io/docs/model/rforest.html} for an example in Radiant
\end{Details}
%
\begin{SeeAlso}
\code{\LinkA{rforest}{rforest}} to generate the result

\code{\LinkA{summary.rforest}{summary.rforest}} to summarize results
\end{SeeAlso}
%
\begin{Examples}
\begin{ExampleCode}
result <- rforest(titanic, "survived", c("pclass", "sex"), lev = "Yes")
predict(result, pred_cmd = "pclass = levels(pclass)")
result <- rforest(diamonds, "price", "carat:color", type = "regression")
predict(result, pred_cmd = "carat = 1:3")
predict(result, pred_data = diamonds) %>% head()

\end{ExampleCode}
\end{Examples}
\HeaderA{predict\_model}{Predict method for model functions}{predict.Rul.model}
%
\begin{Description}
Predict method for model functions
\end{Description}
%
\begin{Usage}
\begin{verbatim}
predict_model(
  object,
  pfun,
  mclass,
  pred_data = NULL,
  pred_cmd = "",
  conf_lev = 0.95,
  se = FALSE,
  dec = 3,
  envir = parent.frame(),
  ...
)
\end{verbatim}
\end{Usage}
%
\begin{Arguments}
\begin{ldescription}
\item[\code{object}] Return value from \code{\LinkA{regress}{regress}}

\item[\code{pfun}] Function to use for prediction

\item[\code{mclass}] Model class to attach

\item[\code{pred\_data}] Dataset to use for prediction

\item[\code{pred\_cmd}] Command used to generate data for prediction (e.g., 'carat = 1:10')

\item[\code{conf\_lev}] Confidence level used to estimate confidence intervals (.95 is the default)

\item[\code{se}] Logical that indicates if prediction standard errors should be calculated (default = FALSE)

\item[\code{dec}] Number of decimals to show

\item[\code{envir}] Environment to extract data from

\item[\code{...}] further arguments passed to or from other methods
\end{ldescription}
\end{Arguments}
%
\begin{Details}
See \url{https://radiant-rstats.github.io/docs/model/regress.html} for an example in Radiant
\end{Details}
\HeaderA{pred\_plot}{Prediction Plots}{pred.Rul.plot}
%
\begin{Description}
Prediction Plots
\end{Description}
%
\begin{Usage}
\begin{verbatim}
pred_plot(
  x,
  plot_list = list(),
  incl,
  incl_int,
  fix = TRUE,
  hline = TRUE,
  nr = 20,
  pdp_range = c(0.025, 0.975),
  minq = NULL,
  maxq = NULL
)
\end{verbatim}
\end{Usage}
%
\begin{Arguments}
\begin{ldescription}
\item[\code{x}] Return value from a model

\item[\code{plot\_list}] List used to store plots

\item[\code{incl}] Which variables to include in prediction plots

\item[\code{incl\_int}] Which interactions to investigate in prediction plots

\item[\code{fix}] Set the desired limited on yhat or have it calculated automatically.
Set to FALSE to have y-axis limits set by ggplot2 for each plot

\item[\code{hline}] Add a dashed horizontal line at the mean response. Set to FALSE to suppress,
or a numeric value to draw the line at that specific value

\item[\code{nr}] Number of values to use to generate predictions for a numeric explanatory variable

\item[\code{pdp\_range}] Numeric vector \code{c(lo, hi)} giving the percentile range used to trim the
x-axis for numeric predictors (default \code{c(0.025, 0.975)})

\item[\code{minq}] Deprecated. Use \code{pdp\_range[1]} instead

\item[\code{maxq}] Deprecated. Use \code{pdp\_range[2]} instead
\end{ldescription}
\end{Arguments}
%
\begin{Details}
Faster, but less robust, alternative for PDP plots. Variable
values not included in the prediction are set to either the mean or
the most common value (level)
\end{Details}
\HeaderA{print.crtree.predict}{Print method for predict.crtree}{print.crtree.predict}
%
\begin{Description}
Print method for predict.crtree
\end{Description}
%
\begin{Usage}
\begin{verbatim}
## S3 method for class 'crtree.predict'
print(x, ..., n = 10)
\end{verbatim}
\end{Usage}
%
\begin{Arguments}
\begin{ldescription}
\item[\code{x}] Return value from prediction method

\item[\code{...}] further arguments passed to or from other methods

\item[\code{n}] Number of lines of prediction results to print. Use -1 to print all lines
\end{ldescription}
\end{Arguments}
\HeaderA{print.gbt.predict}{Print method for predict.gbt}{print.gbt.predict}
%
\begin{Description}
Print method for predict.gbt
\end{Description}
%
\begin{Usage}
\begin{verbatim}
## S3 method for class 'gbt.predict'
print(x, ..., n = 10)
\end{verbatim}
\end{Usage}
%
\begin{Arguments}
\begin{ldescription}
\item[\code{x}] Return value from prediction method

\item[\code{...}] further arguments passed to or from other methods

\item[\code{n}] Number of lines of prediction results to print. Use -1 to print all lines
\end{ldescription}
\end{Arguments}
\HeaderA{print.logistic.predict}{Print method for logistic.predict}{print.logistic.predict}
%
\begin{Description}
Print method for logistic.predict
\end{Description}
%
\begin{Usage}
\begin{verbatim}
## S3 method for class 'logistic.predict'
print(x, ..., n = 10)
\end{verbatim}
\end{Usage}
%
\begin{Arguments}
\begin{ldescription}
\item[\code{x}] Return value from prediction method

\item[\code{...}] further arguments passed to or from other methods

\item[\code{n}] Number of lines of prediction results to print. Use -1 to print all lines
\end{ldescription}
\end{Arguments}
\HeaderA{print.mnl.predict}{Print method for mnl.predict}{print.mnl.predict}
%
\begin{Description}
Print method for mnl.predict
\end{Description}
%
\begin{Usage}
\begin{verbatim}
## S3 method for class 'mnl.predict'
print(x, ..., n = 10)
\end{verbatim}
\end{Usage}
%
\begin{Arguments}
\begin{ldescription}
\item[\code{x}] Return value from prediction method

\item[\code{...}] further arguments passed to or from other methods

\item[\code{n}] Number of lines of prediction results to print. Use -1 to print all lines
\end{ldescription}
\end{Arguments}
\HeaderA{print.nb.predict}{Print method for predict.nb}{print.nb.predict}
%
\begin{Description}
Print method for predict.nb
\end{Description}
%
\begin{Usage}
\begin{verbatim}
## S3 method for class 'nb.predict'
print(x, ..., n = 10)
\end{verbatim}
\end{Usage}
%
\begin{Arguments}
\begin{ldescription}
\item[\code{x}] Return value from prediction method

\item[\code{...}] further arguments passed to or from other methods

\item[\code{n}] Number of lines of prediction results to print. Use -1 to print all lines
\end{ldescription}
\end{Arguments}
\HeaderA{print.nn.predict}{Print method for predict.nn}{print.nn.predict}
%
\begin{Description}
Print method for predict.nn
\end{Description}
%
\begin{Usage}
\begin{verbatim}
## S3 method for class 'nn.predict'
print(x, ..., n = 10)
\end{verbatim}
\end{Usage}
%
\begin{Arguments}
\begin{ldescription}
\item[\code{x}] Return value from prediction method

\item[\code{...}] further arguments passed to or from other methods

\item[\code{n}] Number of lines of prediction results to print. Use -1 to print all lines
\end{ldescription}
\end{Arguments}
\HeaderA{print.regress.predict}{Print method for predict.regress}{print.regress.predict}
%
\begin{Description}
Print method for predict.regress
\end{Description}
%
\begin{Usage}
\begin{verbatim}
## S3 method for class 'regress.predict'
print(x, ..., n = 10)
\end{verbatim}
\end{Usage}
%
\begin{Arguments}
\begin{ldescription}
\item[\code{x}] Return value from prediction method

\item[\code{...}] further arguments passed to or from other methods

\item[\code{n}] Number of lines of prediction results to print. Use -1 to print all lines
\end{ldescription}
\end{Arguments}
\HeaderA{print.rforest.predict}{Print method for predict.rforest}{print.rforest.predict}
%
\begin{Description}
Print method for predict.rforest
\end{Description}
%
\begin{Usage}
\begin{verbatim}
## S3 method for class 'rforest.predict'
print(x, ..., n = 10)
\end{verbatim}
\end{Usage}
%
\begin{Arguments}
\begin{ldescription}
\item[\code{x}] Return value from prediction method

\item[\code{...}] further arguments passed to or from other methods

\item[\code{n}] Number of lines of prediction results to print. Use -1 to print all lines
\end{ldescription}
\end{Arguments}
\HeaderA{print\_predict\_model}{Print method for the model prediction}{print.Rul.predict.Rul.model}
%
\begin{Description}
Print method for the model prediction
\end{Description}
%
\begin{Usage}
\begin{verbatim}
print_predict_model(x, ..., n = 10, header = "")
\end{verbatim}
\end{Usage}
%
\begin{Arguments}
\begin{ldescription}
\item[\code{x}] Return value from prediction method

\item[\code{...}] further arguments passed to or from other methods

\item[\code{n}] Number of lines of prediction results to print. Use -1 to print all lines

\item[\code{header}] Header line
\end{ldescription}
\end{Arguments}
\HeaderA{profit}{Calculate Profit based on cost:margin ratio}{profit}
%
\begin{Description}
Calculate Profit based on cost:margin ratio
\end{Description}
%
\begin{Usage}
\begin{verbatim}
profit(pred, rvar, lev, cost = 1, margin = 2)
\end{verbatim}
\end{Usage}
%
\begin{Arguments}
\begin{ldescription}
\item[\code{pred}] Prediction or predictor

\item[\code{rvar}] Response variable

\item[\code{lev}] The level in the response variable defined as success

\item[\code{cost}] Cost per treatment (e.g., mailing costs)

\item[\code{margin}] Margin, or benefit, per 'success' (e.g., customer purchase). A cost:margin ratio of 1:2 implies
the cost of False Positive are equivalent to the benefits of a True Positive
\end{ldescription}
\end{Arguments}
%
\begin{Value}
profit
\end{Value}
%
\begin{Examples}
\begin{ExampleCode}
profit(runif(20000), dvd$buy, "yes", cost = 1, margin = 2)
profit(ifelse(dvd$buy == "yes", 1, 0), dvd$buy, "yes", cost = 1, margin = 20)
profit(ifelse(dvd$buy == "yes", 1, 0), dvd$buy)
\end{ExampleCode}
\end{Examples}
\HeaderA{radiant.model}{radiant.model}{radiant.model}
%
\begin{Description}
Launch radiant.model in the default web browser
\end{Description}
%
\begin{Usage}
\begin{verbatim}
radiant.model(state, ...)
\end{verbatim}
\end{Usage}
%
\begin{Arguments}
\begin{ldescription}
\item[\code{state}] Path to state file to load

\item[\code{...}] additional arguments to pass to shiny::runApp (e.g, port = 8080)
\end{ldescription}
\end{Arguments}
%
\begin{Details}
See \url{https://radiant-rstats.github.io/docs/} for documentation and tutorials
\end{Details}
%
\begin{Examples}
\begin{ExampleCode}
## Not run: 
radiant.model()

## End(Not run)
\end{ExampleCode}
\end{Examples}
\HeaderA{radiant.model-deprecated}{Deprecated function(s) in the radiant.model package}{radiant.model.Rdash.deprecated}
\aliasA{ann}{radiant.model-deprecated}{ann}
%
\begin{Description}
These functions are provided for compatibility with previous versions of
radiant. They will eventually be  removed.
\end{Description}
%
\begin{Usage}
\begin{verbatim}
ann(...)
\end{verbatim}
\end{Usage}
%
\begin{Arguments}
\begin{ldescription}
\item[\code{...}] Parameters to be passed to the updated functions
\end{ldescription}
\end{Arguments}
%
\begin{Section}{Details}


\Tabular{rl}{
\code{ann} is now a synonym for \code{\LinkA{nn}{nn}}\\{}
\code{scaledf} is now a synonym for \code{\LinkA{scale\_df}{scale.Rul.df}}\\{}
}
\end{Section}
\HeaderA{radiant.model\_viewer}{Launch radiant.model in the Rstudio viewer}{radiant.model.Rul.viewer}
%
\begin{Description}
Launch radiant.model in the Rstudio viewer
\end{Description}
%
\begin{Usage}
\begin{verbatim}
radiant.model_viewer(state, ...)
\end{verbatim}
\end{Usage}
%
\begin{Arguments}
\begin{ldescription}
\item[\code{state}] Path to state file to load

\item[\code{...}] additional arguments to pass to shiny::runApp (e.g, port = 8080)
\end{ldescription}
\end{Arguments}
%
\begin{Details}
See \url{https://radiant-rstats.github.io/docs/} for documentation and tutorials
\end{Details}
%
\begin{Examples}
\begin{ExampleCode}
## Not run: 
radiant.model_viewer()

## End(Not run)
\end{ExampleCode}
\end{Examples}
\HeaderA{radiant.model\_window}{Launch radiant.model in an Rstudio window}{radiant.model.Rul.window}
%
\begin{Description}
Launch radiant.model in an Rstudio window
\end{Description}
%
\begin{Usage}
\begin{verbatim}
radiant.model_window(state, ...)
\end{verbatim}
\end{Usage}
%
\begin{Arguments}
\begin{ldescription}
\item[\code{state}] Path to state file to load

\item[\code{...}] additional arguments to pass to shiny::runApp (e.g, port = 8080)
\end{ldescription}
\end{Arguments}
%
\begin{Details}
See \url{https://radiant-rstats.github.io/docs/} for documentation and tutorials
\end{Details}
%
\begin{Examples}
\begin{ExampleCode}
## Not run: 
radiant.model_window()

## End(Not run)
\end{ExampleCode}
\end{Examples}
\HeaderA{ratings}{Movie ratings}{ratings}
\keyword{datasets}{ratings}
%
\begin{Description}
Movie ratings
\end{Description}
%
\begin{Usage}
\begin{verbatim}
data(ratings)
\end{verbatim}
\end{Usage}
%
\begin{Format}
A data frame with 110 rows and 4 variables
\end{Format}
%
\begin{Details}
Use collaborative filtering to create recommendations based on ratings from existing users. Description provided in attr(ratings, "description")
\end{Details}
\HeaderA{regress}{Linear regression using OLS}{regress}
%
\begin{Description}
Linear regression using OLS
\end{Description}
%
\begin{Usage}
\begin{verbatim}
regress(
  dataset,
  rvar,
  evar,
  int = "",
  check = "",
  form,
  data_filter = "",
  arr = "",
  rows = NULL,
  envir = parent.frame()
)
\end{verbatim}
\end{Usage}
%
\begin{Arguments}
\begin{ldescription}
\item[\code{dataset}] Dataset

\item[\code{rvar}] The response variable in the regression

\item[\code{evar}] Explanatory variables in the regression

\item[\code{int}] Interaction terms to include in the model

\item[\code{check}] Use "standardize" to see standardized coefficient estimates. Use "stepwise-backward" (or "stepwise-forward", or "stepwise-both") to apply step-wise selection of variables in estimation. Add "robust" for robust estimation of standard errors (HC1)

\item[\code{form}] Optional formula to use instead of rvar, evar, and int

\item[\code{data\_filter}] Expression entered in, e.g., Data > View to filter the dataset in Radiant. The expression should be a string (e.g., "price > 10000")

\item[\code{arr}] Expression to arrange (sort) the data on (e.g., "color, desc(price)")

\item[\code{rows}] Rows to select from the specified dataset

\item[\code{envir}] Environment to extract data from
\end{ldescription}
\end{Arguments}
%
\begin{Details}
See \url{https://radiant-rstats.github.io/docs/model/regress.html} for an example in Radiant
\end{Details}
%
\begin{Value}
A list of all variables used in the regress function as an object of class regress
\end{Value}
%
\begin{SeeAlso}
\code{\LinkA{summary.regress}{summary.regress}} to summarize results

\code{\LinkA{plot.regress}{plot.regress}} to plot results

\code{\LinkA{predict.regress}{predict.regress}} to generate predictions
\end{SeeAlso}
%
\begin{Examples}
\begin{ExampleCode}
regress(diamonds, "price", c("carat", "clarity"), check = "standardize") %>% summary()
regress(diamonds, "price", c("carat", "clarity")) %>% str()

\end{ExampleCode}
\end{Examples}
\HeaderA{remove\_comments}{Remove comments from formula before it is evaluated}{remove.Rul.comments}
%
\begin{Description}
Remove comments from formula before it is evaluated
\end{Description}
%
\begin{Usage}
\begin{verbatim}
remove_comments(x)
\end{verbatim}
\end{Usage}
%
\begin{Arguments}
\begin{ldescription}
\item[\code{x}] Input string
\end{ldescription}
\end{Arguments}
%
\begin{Value}
Cleaned string
\end{Value}
\HeaderA{render.DiagrammeR}{Method to render DiagrammeR plots}{render.DiagrammeR}
%
\begin{Description}
Method to render DiagrammeR plots
\end{Description}
%
\begin{Usage}
\begin{verbatim}
## S3 method for class 'DiagrammeR'
render(object, shiny = shiny::getDefaultReactiveDomain(), ...)
\end{verbatim}
\end{Usage}
%
\begin{Arguments}
\begin{ldescription}
\item[\code{object}] DiagrammeR plot

\item[\code{shiny}] Check if function is called from a shiny application

\item[\code{...}] Additional arguments
\end{ldescription}
\end{Arguments}
\HeaderA{repeater}{Repeated simulation}{repeater}
%
\begin{Description}
Repeated simulation
\end{Description}
%
\begin{Usage}
\begin{verbatim}
repeater(
  dataset,
  nr = 12,
  vars = "",
  grid = "",
  sum_vars = "",
  byvar = ".sim",
  fun = "sum",
  form = "",
  seed = NULL,
  name = "",
  envir = parent.frame()
)
\end{verbatim}
\end{Usage}
%
\begin{Arguments}
\begin{ldescription}
\item[\code{dataset}] Return value from the simulater function

\item[\code{nr}] Number times to repeat the simulation

\item[\code{vars}] Variables to use in repeated simulation

\item[\code{grid}] Character vector of expressions to use in grid search for constants

\item[\code{sum\_vars}] (Numeric) variables to summaries

\item[\code{byvar}] Variable(s) to group data by before summarizing

\item[\code{fun}] Functions to use for summarizing

\item[\code{form}] A character vector with the formula to apply to the summarized data

\item[\code{seed}] Seed for the repeated simulation

\item[\code{name}] Deprecated argument

\item[\code{envir}] Environment to extract data from
\end{ldescription}
\end{Arguments}
%
\begin{SeeAlso}
\code{\LinkA{summary.repeater}{summary.repeater}} to summarize results from repeated simulation

\code{\LinkA{plot.repeater}{plot.repeater}} to plot results from repeated simulation
\end{SeeAlso}
%
\begin{Examples}
\begin{ExampleCode}
simdat <- simulater(
  const = c("var_cost 5", "fixed_cost 1000"),
  norm = "E 0 100;",
  discrete = "price 6 8 .3 .7;",
  form = c(
    "demand = 1000 - 50*price + E",
    "profit = demand*(price-var_cost) - fixed_cost",
    "profit_small = profit < 100"
  ),
  seed = 1234
)

repdat <- repeater(
  simdat,
  nr = 12,
  vars = c("E", "price"),
  sum_vars = "profit",
  byvar = ".sim",
  form = "profit_365 = profit_sum < 36500",
  seed = 1234,
)

head(repdat)
summary(repdat)
plot(repdat)

\end{ExampleCode}
\end{Examples}
\HeaderA{rforest}{Random Forest using Ranger}{rforest}
%
\begin{Description}
Random Forest using Ranger
\end{Description}
%
\begin{Usage}
\begin{verbatim}
rforest(
  dataset,
  rvar,
  evar,
  type = "classification",
  lev = "",
  mtry = NULL,
  num.trees = 100,
  min.node.size = 1,
  sample.fraction = 1,
  replace = NULL,
  num.threads = 12,
  wts = "None",
  seed = NA,
  data_filter = "",
  arr = "",
  rows = NULL,
  envir = parent.frame(),
  ...
)
\end{verbatim}
\end{Usage}
%
\begin{Arguments}
\begin{ldescription}
\item[\code{dataset}] Dataset

\item[\code{rvar}] The response variable in the model

\item[\code{evar}] Explanatory variables in the model

\item[\code{type}] Model type (i.e., "classification" or "regression")

\item[\code{lev}] Level to use as the first column in prediction output

\item[\code{mtry}] Number of variables to possibly split at in each node. Default is the (rounded down) square root of the number variables

\item[\code{num.trees}] Number of trees to create

\item[\code{min.node.size}] Minimal node size

\item[\code{sample.fraction}] Fraction of observations to sample. Default is 1 for sampling with replacement and 0.632 for sampling without replacement

\item[\code{replace}] Sample with (TRUE) or without (FALSE) replacement. If replace is NULL it will be reset to TRUE if the sample.fraction is equal to 1 and will be set to FALSE otherwise

\item[\code{num.threads}] Number of parallel threads to use. Defaults to 12 if available

\item[\code{wts}] Case weights to use in estimation

\item[\code{seed}] Random seed to use as the starting point

\item[\code{data\_filter}] Expression entered in, e.g., Data > View to filter the dataset in Radiant. The expression should be a string (e.g., "price > 10000")

\item[\code{arr}] Expression to arrange (sort) the data on (e.g., "color, desc(price)")

\item[\code{rows}] Rows to select from the specified dataset

\item[\code{envir}] Environment to extract data from

\item[\code{...}] Further arguments to pass to ranger
\end{ldescription}
\end{Arguments}
%
\begin{Details}
See \url{https://radiant-rstats.github.io/docs/model/rforest.html} for an example in Radiant
\end{Details}
%
\begin{Value}
A list with all variables defined in rforest as an object of class rforest
\end{Value}
%
\begin{SeeAlso}
\code{\LinkA{summary.rforest}{summary.rforest}} to summarize results

\code{\LinkA{plot.rforest}{plot.rforest}} to plot results

\code{\LinkA{predict.rforest}{predict.rforest}} for prediction
\end{SeeAlso}
%
\begin{Examples}
\begin{ExampleCode}
rforest(titanic, "survived", c("pclass", "sex"), lev = "Yes") %>% summary()
rforest(titanic, "survived", c("pclass", "sex")) %>% str()
rforest(titanic, "survived", c("pclass", "sex"), max.depth = 1)
rforest(diamonds, "price", c("carat", "clarity"), type = "regression") %>% summary()

\end{ExampleCode}
\end{Examples}
\HeaderA{rig}{Relative Information Gain (RIG)}{rig}
%
\begin{Description}
Relative Information Gain (RIG)
\end{Description}
%
\begin{Usage}
\begin{verbatim}
rig(pred, rvar, lev, crv = 1e-07, na.rm = TRUE)
\end{verbatim}
\end{Usage}
%
\begin{Arguments}
\begin{ldescription}
\item[\code{pred}] Prediction or predictor

\item[\code{rvar}] Response variable

\item[\code{lev}] The level in the response variable defined as success

\item[\code{crv}] Correction value to avoid log(0)

\item[\code{na.rm}] Logical that indicates if missing values should be removed (TRUE) or not (FALSE)
\end{ldescription}
\end{Arguments}
%
\begin{Details}
See \url{https://radiant-rstats.github.io/docs/model/evalbin.html} for an example in Radiant
\end{Details}
%
\begin{Value}
RIG statistic
\end{Value}
%
\begin{SeeAlso}
\code{\LinkA{evalbin}{evalbin}} to calculate results

\code{\LinkA{summary.evalbin}{summary.evalbin}} to summarize results

\code{\LinkA{plot.evalbin}{plot.evalbin}} to plot results
\end{SeeAlso}
%
\begin{Examples}
\begin{ExampleCode}
rig(runif(20000), dvd$buy, "yes")
rig(ifelse(dvd$buy == "yes", 1, 0), dvd$buy, "yes")
\end{ExampleCode}
\end{Examples}
\HeaderA{RMSE}{Root Mean Squared Error}{RMSE}
%
\begin{Description}
Root Mean Squared Error
\end{Description}
%
\begin{Usage}
\begin{verbatim}
RMSE(pred, rvar)
\end{verbatim}
\end{Usage}
%
\begin{Arguments}
\begin{ldescription}
\item[\code{pred}] Prediction (vector)

\item[\code{rvar}] Response (vector)
\end{ldescription}
\end{Arguments}
%
\begin{Value}
Root Mean Squared Error
\end{Value}
\HeaderA{Rsq}{R-squared}{Rsq}
%
\begin{Description}
R-squared
\end{Description}
%
\begin{Usage}
\begin{verbatim}
Rsq(pred, rvar)
\end{verbatim}
\end{Usage}
%
\begin{Arguments}
\begin{ldescription}
\item[\code{pred}] Prediction (vector)

\item[\code{rvar}] Response (vector)
\end{ldescription}
\end{Arguments}
%
\begin{Value}
R-squared
\end{Value}
\HeaderA{scale\_df}{Center or standardize variables in a data frame}{scale.Rul.df}
%
\begin{Description}
Center or standardize variables in a data frame
\end{Description}
%
\begin{Usage}
\begin{verbatim}
scale_df(dataset, center = TRUE, scale = TRUE, sf = 2, wts = NULL, calc = TRUE)
\end{verbatim}
\end{Usage}
%
\begin{Arguments}
\begin{ldescription}
\item[\code{dataset}] Data frame

\item[\code{center}] Center data (TRUE or FALSE)

\item[\code{scale}] Scale data (TRUE or FALSE)

\item[\code{sf}] Scaling factor (default is 2)

\item[\code{wts}] Weights to use (default is NULL for no weights)

\item[\code{calc}] Calculate mean and sd or use attributes attached to dat
\end{ldescription}
\end{Arguments}
%
\begin{Value}
Scaled data frame
\end{Value}
\HeaderA{sdw}{Standard deviation of weighted sum of variables}{sdw}
%
\begin{Description}
Standard deviation of weighted sum of variables
\end{Description}
%
\begin{Usage}
\begin{verbatim}
sdw(...)
\end{verbatim}
\end{Usage}
%
\begin{Arguments}
\begin{ldescription}
\item[\code{...}] A matched number of weights and stocks
\end{ldescription}
\end{Arguments}
%
\begin{Value}
A vector of standard deviation estimates
\end{Value}
\HeaderA{sensitivity}{Method to evaluate sensitivity of an analysis}{sensitivity}
%
\begin{Description}
Method to evaluate sensitivity of an analysis
\end{Description}
%
\begin{Usage}
\begin{verbatim}
sensitivity(object, ...)
\end{verbatim}
\end{Usage}
%
\begin{Arguments}
\begin{ldescription}
\item[\code{object}] Object of relevant class for which to evaluate sensitivity

\item[\code{...}] Additional arguments
\end{ldescription}
\end{Arguments}
%
\begin{SeeAlso}
\code{\LinkA{sensitivity.dtree}{sensitivity.dtree}} to plot results
\end{SeeAlso}
\HeaderA{sensitivity.dtree}{Evaluate sensitivity of the decision tree}{sensitivity.dtree}
%
\begin{Description}
Evaluate sensitivity of the decision tree
\end{Description}
%
\begin{Usage}
\begin{verbatim}
## S3 method for class 'dtree'
sensitivity(
  object,
  vars = NULL,
  decs = NULL,
  envir = parent.frame(),
  shiny = FALSE,
  custom = FALSE,
  ...
)
\end{verbatim}
\end{Usage}
%
\begin{Arguments}
\begin{ldescription}
\item[\code{object}] Return value from \code{\LinkA{dtree}{dtree}}

\item[\code{vars}] Variables to include in the sensitivity analysis

\item[\code{decs}] Decisions to include in the sensitivity analysis

\item[\code{envir}] Environment to extract data from

\item[\code{shiny}] Did the function call originate inside a shiny app

\item[\code{custom}] Logical (TRUE, FALSE) to indicate if ggplot object (or list of ggplot objects) should be returned. This option can be used to customize plots (e.g., add a title, change x and y labels, etc.). See examples and \url{https://ggplot2.tidyverse.org} for options.

\item[\code{...}] Additional arguments
\end{ldescription}
\end{Arguments}
%
\begin{Details}
See \url{https://radiant-rstats.github.io/docs/model/dtree.html} for an example in Radiant
\end{Details}
%
\begin{SeeAlso}
\code{\LinkA{dtree}{dtree}} to generate the result

\code{\LinkA{plot.dtree}{plot.dtree}} to summarize results

\code{\LinkA{summary.dtree}{summary.dtree}} to summarize results
\end{SeeAlso}
%
\begin{Examples}
\begin{ExampleCode}
dtree(movie_contract, opt = "max") %>%
  sensitivity(
    vars = "legal fees 0 100000 10000",
    decs = c("Sign with Movie Company", "Sign with TV Network"),
    custom = FALSE
  )

\end{ExampleCode}
\end{Examples}
\HeaderA{simulater}{Simulate data for decision analysis}{simulater}
%
\begin{Description}
Simulate data for decision analysis
\end{Description}
%
\begin{Usage}
\begin{verbatim}
simulater(
  const = "",
  lnorm = "",
  norm = "",
  unif = "",
  discrete = "",
  binom = "",
  pois = "",
  sequ = "",
  grid = "",
  data = NULL,
  form = "",
  funcs = "",
  seed = NULL,
  nexact = FALSE,
  ncorr = NULL,
  name = "",
  nr = 1000,
  dataset = NULL,
  envir = parent.frame()
)
\end{verbatim}
\end{Usage}
%
\begin{Arguments}
\begin{ldescription}
\item[\code{const}] A character vector listing the constants to include in the analysis (e.g., c("cost = 3", "size = 4"))

\item[\code{lnorm}] A character vector listing the log-normally distributed random variables to include in the analysis (e.g., "demand 2000 1000" where the first number is the log-mean and the second is the log-standard deviation)

\item[\code{norm}] A character vector listing the normally distributed random variables to include in the analysis (e.g., "demand 2000 1000" where the first number is the mean and the second is the standard deviation)

\item[\code{unif}] A character vector listing the uniformly distributed random variables to include in the analysis (e.g., "demand 0 1" where the first number is the minimum value and the second is the maximum value)

\item[\code{discrete}] A character vector listing the random variables with a discrete distribution to include in the analysis (e.g., "price 5 8 .3 .7" where the first set of numbers are the values and the second set the probabilities

\item[\code{binom}] A character vector listing the random variables with a binomial distribution to include in the analysis (e.g., "crash 100 .01") where the first number is the number of trials and the second is the probability of success)

\item[\code{pois}] A character vector listing the random variables with a poisson distribution to include in the analysis (e.g., "demand 10") where the number is the lambda value (i.e., the average number of events or the event rate)

\item[\code{sequ}] A character vector listing the start and end for a sequence to include in the analysis (e.g., "trend 1 100 1"). The number of 'steps' is determined by the number of simulations

\item[\code{grid}] A character vector listing the start, end, and step for a set of sequences to include in the analysis (e.g., "trend 1 100 1"). The number of rows in the expanded will over ride the number of simulations

\item[\code{data}] Dataset to be used in the calculations

\item[\code{form}] A character vector with the formula to evaluate (e.g., "profit = demand * (price - cost)")

\item[\code{funcs}] A named list of user defined functions to apply to variables generated as part of the simulation

\item[\code{seed}] Optional seed used in simulation

\item[\code{nexact}] Logical to indicate if normally distributed random variables should be simulated to the exact specified values

\item[\code{ncorr}] A string of correlations used for normally distributed random variables. The number of values should be equal to one or to the number of combinations of variables simulated

\item[\code{name}] Deprecated argument

\item[\code{nr}] Number of simulations

\item[\code{dataset}] Data list from previous simulation. Used by repeater function

\item[\code{envir}] Environment to extract data from
\end{ldescription}
\end{Arguments}
%
\begin{Details}
See \url{https://radiant-rstats.github.io/docs/model/simulater.html} for an example in Radiant
\end{Details}
%
\begin{Value}
A data.frame with the simulated data
\end{Value}
%
\begin{SeeAlso}
\code{\LinkA{summary.simulater}{summary.simulater}} to summarize results

\code{\LinkA{plot.simulater}{plot.simulater}} to plot results
\end{SeeAlso}
%
\begin{Examples}
\begin{ExampleCode}
simulater(
  const = "cost 3",
  norm = "demand 2000 1000",
  discrete = "price 5 8 .3 .7",
  form = "profit = demand * (price - cost)",
  seed = 1234
) %>% str()

\end{ExampleCode}
\end{Examples}
\HeaderA{sim\_cleaner}{Clean input command string}{sim.Rul.cleaner}
%
\begin{Description}
Clean input command string
\end{Description}
%
\begin{Usage}
\begin{verbatim}
sim_cleaner(x)
\end{verbatim}
\end{Usage}
%
\begin{Arguments}
\begin{ldescription}
\item[\code{x}] Input string
\end{ldescription}
\end{Arguments}
%
\begin{Value}
Cleaned string
\end{Value}
\HeaderA{sim\_cor}{Simulate correlated normally distributed data}{sim.Rul.cor}
%
\begin{Description}
Simulate correlated normally distributed data
\end{Description}
%
\begin{Usage}
\begin{verbatim}
sim_cor(n, rho, means, sds, exact = FALSE)
\end{verbatim}
\end{Usage}
%
\begin{Arguments}
\begin{ldescription}
\item[\code{n}] The number of values to simulate (i.e., the number of rows in the simulated data)

\item[\code{rho}] A vector of correlations to apply to the columns of the simulated data. The number of values should be equal to one or to the number of combinations of variables to be simulated

\item[\code{means}] A vector of means. The number of values should be equal to the number of variables to simulate

\item[\code{sds}] A vector of standard deviations. The number of values should be equal to the number of variables to simulate

\item[\code{exact}] A logical that indicates if the inputs should be interpreted as population of sample characteristics
\end{ldescription}
\end{Arguments}
%
\begin{Value}
A data.frame with the simulated data
\end{Value}
%
\begin{Examples}
\begin{ExampleCode}
sim <- sim_cor(100, .74, c(0, 10), c(1, 5), exact = TRUE)
cor(sim)
sim_summary(sim)

\end{ExampleCode}
\end{Examples}
\HeaderA{sim\_splitter}{Split input command string}{sim.Rul.splitter}
%
\begin{Description}
Split input command string
\end{Description}
%
\begin{Usage}
\begin{verbatim}
sim_splitter(x, symbol = " ")
\end{verbatim}
\end{Usage}
%
\begin{Arguments}
\begin{ldescription}
\item[\code{x}] Input string

\item[\code{symbol}] Symbol used to split the command string
\end{ldescription}
\end{Arguments}
%
\begin{Value}
Split input command string
\end{Value}
\HeaderA{sim\_summary}{Print simulation summary}{sim.Rul.summary}
%
\begin{Description}
Print simulation summary
\end{Description}
%
\begin{Usage}
\begin{verbatim}
sim_summary(dataset, dc = get_class(dataset), fun = "", dec = 4)
\end{verbatim}
\end{Usage}
%
\begin{Arguments}
\begin{ldescription}
\item[\code{dataset}] Simulated data

\item[\code{dc}] Variable classes

\item[\code{fun}] Summary function to apply

\item[\code{dec}] Number of decimals to show
\end{ldescription}
\end{Arguments}
%
\begin{SeeAlso}
\code{\LinkA{simulater}{simulater}} to run a simulation

\code{\LinkA{repeater}{repeater}} to run a repeated simulation
\end{SeeAlso}
%
\begin{Examples}
\begin{ExampleCode}
simulater(
  const = "cost 3",
  norm = "demand 2000 1000",
  discrete = "price 5 8 .3 .7",
  form = c("profit = demand * (price - cost)", "profit5K = profit > 5000"),
  seed = 1234
) %>% sim_summary()

\end{ExampleCode}
\end{Examples}
\HeaderA{store.crs}{Deprecated: Store method for the crs function}{store.crs}
%
\begin{Description}
Deprecated: Store method for the crs function
\end{Description}
%
\begin{Usage}
\begin{verbatim}
## S3 method for class 'crs'
store(dataset, object, name, ...)
\end{verbatim}
\end{Usage}
%
\begin{Arguments}
\begin{ldescription}
\item[\code{dataset}] Dataset

\item[\code{object}] Return value from \code{\LinkA{crs}{crs}}

\item[\code{name}] Name to assign to the dataset

\item[\code{...}] further arguments passed to or from other methods
\end{ldescription}
\end{Arguments}
%
\begin{Details}
Return recommendations See \url{https://radiant-rstats.github.io/docs/model/crs.html} for an example in Radiant
\end{Details}
\HeaderA{store.mnl.predict}{Store predicted values generated in the mnl function}{store.mnl.predict}
%
\begin{Description}
Store predicted values generated in the mnl function
\end{Description}
%
\begin{Usage}
\begin{verbatim}
## S3 method for class 'mnl.predict'
store(dataset, object, name = NULL, ...)
\end{verbatim}
\end{Usage}
%
\begin{Arguments}
\begin{ldescription}
\item[\code{dataset}] Dataset to add predictions to

\item[\code{object}] Return value from model function

\item[\code{name}] Variable name(s) assigned to predicted values. If empty, the levels of the response variable will be used

\item[\code{...}] Additional arguments
\end{ldescription}
\end{Arguments}
%
\begin{Details}
See \url{https://radiant-rstats.github.io/docs/model/mnl.html} for an example in Radiant
\end{Details}
%
\begin{Examples}
\begin{ExampleCode}
result <- mnl(
  ketchup,
  rvar = "choice",
  evar = c("price.heinz28", "price.heinz32", "price.heinz41", "price.hunts32"),
  lev = "heinz28"
)
pred <- predict(result, pred_data = ketchup)
ketchup <- store(ketchup, pred, name = c("heinz28", "heinz32", "heinz41", "hunts32"))

\end{ExampleCode}
\end{Examples}
\HeaderA{store.model}{Store residuals from a model}{store.model}
%
\begin{Description}
Store residuals from a model
\end{Description}
%
\begin{Usage}
\begin{verbatim}
## S3 method for class 'model'
store(dataset, object, name = "residuals", ...)
\end{verbatim}
\end{Usage}
%
\begin{Arguments}
\begin{ldescription}
\item[\code{dataset}] Dataset to append residuals to

\item[\code{object}] Return value from a model function

\item[\code{name}] Variable name(s) assigned to model residuals

\item[\code{...}] Additional arguments
\end{ldescription}
\end{Arguments}
%
\begin{Details}
The store method for objects of class "model". Adds model residuals to the dataset while handling missing values and filters. See \url{https://radiant-rstats.github.io/docs/model/regress.html} for an example in Radiant
\end{Details}
%
\begin{Examples}
\begin{ExampleCode}
regress(diamonds, rvar = "price", evar = c("carat", "cut"), data_filter = "price > 1000") %>%
  store(diamonds, ., name = "resid") %>%
  head()

\end{ExampleCode}
\end{Examples}
\HeaderA{store.model.predict}{Store predicted values generated in model functions}{store.model.predict}
%
\begin{Description}
Store predicted values generated in model functions
\end{Description}
%
\begin{Usage}
\begin{verbatim}
## S3 method for class 'model.predict'
store(dataset, object, name = "prediction", ...)
\end{verbatim}
\end{Usage}
%
\begin{Arguments}
\begin{ldescription}
\item[\code{dataset}] Dataset to add predictions to

\item[\code{object}] Return value from model function

\item[\code{name}] Variable name(s) assigned to predicted values

\item[\code{...}] Additional arguments
\end{ldescription}
\end{Arguments}
%
\begin{Details}
See \url{https://radiant-rstats.github.io/docs/model/regress.html} for an example in Radiant
\end{Details}
%
\begin{Examples}
\begin{ExampleCode}
regress(diamonds, rvar = "price", evar = c("carat", "cut")) %>%
  predict(pred_data = diamonds) %>%
  store(diamonds, ., name = c("pred", "pred_low", "pred_high")) %>%
  head()

\end{ExampleCode}
\end{Examples}
\HeaderA{store.nb.predict}{Store predicted values generated in the nb function}{store.nb.predict}
%
\begin{Description}
Store predicted values generated in the nb function
\end{Description}
%
\begin{Usage}
\begin{verbatim}
## S3 method for class 'nb.predict'
store(dataset, object, name = NULL, ...)
\end{verbatim}
\end{Usage}
%
\begin{Arguments}
\begin{ldescription}
\item[\code{dataset}] Dataset to add predictions to

\item[\code{object}] Return value from model function

\item[\code{name}] Variable name(s) assigned to predicted values. If empty, the levels of the response variable will be used

\item[\code{...}] Additional arguments
\end{ldescription}
\end{Arguments}
%
\begin{Details}
See \url{https://radiant-rstats.github.io/docs/model/nb.html} for an example in Radiant
\end{Details}
%
\begin{Examples}
\begin{ExampleCode}
result <- nb(titanic, rvar = "survived", evar = c("pclass", "sex", "age"))
pred <- predict(result, pred_data = titanic)
titanic <- store(titanic, pred, name = c("Yes", "No"))

\end{ExampleCode}
\end{Examples}
\HeaderA{store.rforest.predict}{Store predicted values generated in the rforest function}{store.rforest.predict}
%
\begin{Description}
Store predicted values generated in the rforest function
\end{Description}
%
\begin{Usage}
\begin{verbatim}
## S3 method for class 'rforest.predict'
store(dataset, object, name = NULL, ...)
\end{verbatim}
\end{Usage}
%
\begin{Arguments}
\begin{ldescription}
\item[\code{dataset}] Dataset to add predictions to

\item[\code{object}] Return value from model function

\item[\code{name}] Variable name(s) assigned to predicted values. If empty, the levels of the response variable will be used

\item[\code{...}] Additional arguments
\end{ldescription}
\end{Arguments}
%
\begin{Details}
See \url{https://radiant-rstats.github.io/docs/model/rforest.html} for an example in Radiant
\end{Details}
%
\begin{Examples}
\begin{ExampleCode}
result <- rforest(
  ketchup,
  rvar = "choice",
  evar = c("price.heinz28", "price.heinz32", "price.heinz41", "price.hunts32"),
  lev = "heinz28"
)
pred <- predict(result, pred_data = ketchup)
ketchup <- store(ketchup, pred, name = c("heinz28", "heinz32", "heinz41", "hunts32"))

\end{ExampleCode}
\end{Examples}
\HeaderA{summary.confusion}{Summary method for the confusion matrix}{summary.confusion}
%
\begin{Description}
Summary method for the confusion matrix
\end{Description}
%
\begin{Usage}
\begin{verbatim}
## S3 method for class 'confusion'
summary(object, dec = 3, ...)
\end{verbatim}
\end{Usage}
%
\begin{Arguments}
\begin{ldescription}
\item[\code{object}] Return value from \code{\LinkA{confusion}{confusion}}

\item[\code{dec}] Number of decimals to show

\item[\code{...}] further arguments passed to or from other methods
\end{ldescription}
\end{Arguments}
%
\begin{Details}
See \url{https://radiant-rstats.github.io/docs/model/evalbin.html} for an example in Radiant
\end{Details}
%
\begin{SeeAlso}
\code{\LinkA{confusion}{confusion}} to generate results

\code{\LinkA{plot.confusion}{plot.confusion}} to visualize result
\end{SeeAlso}
%
\begin{Examples}
\begin{ExampleCode}
data.frame(buy = dvd$buy, pred1 = runif(20000), pred2 = ifelse(dvd$buy == "yes", 1, 0)) %>%
  confusion(c("pred1", "pred2"), "buy") %>%
  summary()
\end{ExampleCode}
\end{Examples}
\HeaderA{summary.crs}{Summary method for Collaborative Filter}{summary.crs}
%
\begin{Description}
Summary method for Collaborative Filter
\end{Description}
%
\begin{Usage}
\begin{verbatim}
## S3 method for class 'crs'
summary(object, n = 36, dec = 2, ...)
\end{verbatim}
\end{Usage}
%
\begin{Arguments}
\begin{ldescription}
\item[\code{object}] Return value from \code{\LinkA{crs}{crs}}

\item[\code{n}] Number of lines of recommendations to print. Use -1 to print all lines

\item[\code{dec}] Number of decimals to show

\item[\code{...}] further arguments passed to or from other methods
\end{ldescription}
\end{Arguments}
%
\begin{Details}
See \url{https://radiant-rstats.github.io/docs/model/crs.html} for an example in Radiant
\end{Details}
%
\begin{SeeAlso}
\code{\LinkA{crs}{crs}} to generate the results

\code{\LinkA{plot.crs}{plot.crs}} to plot results if the actual ratings are available
\end{SeeAlso}
%
\begin{Examples}
\begin{ExampleCode}
crs(ratings,
  id = "Users", prod = "Movies", pred = c("M6", "M7", "M8", "M9", "M10"),
  rate = "Ratings", data_filter = "training == 1"
) %>% summary()
\end{ExampleCode}
\end{Examples}
\HeaderA{summary.crtree}{Summary method for the crtree function}{summary.crtree}
%
\begin{Description}
Summary method for the crtree function
\end{Description}
%
\begin{Usage}
\begin{verbatim}
## S3 method for class 'crtree'
summary(object, prn = TRUE, splits = FALSE, cptab = FALSE, modsum = FALSE, ...)
\end{verbatim}
\end{Usage}
%
\begin{Arguments}
\begin{ldescription}
\item[\code{object}] Return value from \code{\LinkA{crtree}{crtree}}

\item[\code{prn}] Print tree in text form

\item[\code{splits}] Print the tree splitting metrics used

\item[\code{cptab}] Print the cp table

\item[\code{modsum}] Print the model summary

\item[\code{...}] further arguments passed to or from other methods
\end{ldescription}
\end{Arguments}
%
\begin{Details}
See \url{https://radiant-rstats.github.io/docs/model/crtree.html} for an example in Radiant
\end{Details}
%
\begin{SeeAlso}
\code{\LinkA{crtree}{crtree}} to generate results

\code{\LinkA{plot.crtree}{plot.crtree}} to plot results

\code{\LinkA{predict.crtree}{predict.crtree}} for prediction
\end{SeeAlso}
%
\begin{Examples}
\begin{ExampleCode}
result <- crtree(titanic, "survived", c("pclass", "sex"), lev = "Yes")
summary(result)
result <- crtree(diamonds, "price", c("carat", "color"), type = "regression")
summary(result)
\end{ExampleCode}
\end{Examples}
\HeaderA{summary.dtree}{Summary method for the dtree function}{summary.dtree}
%
\begin{Description}
Summary method for the dtree function
\end{Description}
%
\begin{Usage}
\begin{verbatim}
## S3 method for class 'dtree'
summary(object, input = TRUE, output = FALSE, dec = 2, ...)
\end{verbatim}
\end{Usage}
%
\begin{Arguments}
\begin{ldescription}
\item[\code{object}] Return value from \code{\LinkA{simulater}{simulater}}

\item[\code{input}] Print decision tree input

\item[\code{output}] Print decision tree output

\item[\code{dec}] Number of decimals to show

\item[\code{...}] further arguments passed to or from other methods
\end{ldescription}
\end{Arguments}
%
\begin{Details}
See \url{https://radiant-rstats.github.io/docs/model/dtree.html} for an example in Radiant
\end{Details}
%
\begin{SeeAlso}
\code{\LinkA{dtree}{dtree}} to generate the results

\code{\LinkA{plot.dtree}{plot.dtree}} to plot results

\code{\LinkA{sensitivity.dtree}{sensitivity.dtree}} to plot results
\end{SeeAlso}
%
\begin{Examples}
\begin{ExampleCode}
dtree(movie_contract, opt = "max") %>% summary(input = TRUE)
dtree(movie_contract, opt = "max") %>% summary(input = FALSE, output = TRUE)

\end{ExampleCode}
\end{Examples}
\HeaderA{summary.evalbin}{Summary method for the evalbin function}{summary.evalbin}
%
\begin{Description}
Summary method for the evalbin function
\end{Description}
%
\begin{Usage}
\begin{verbatim}
## S3 method for class 'evalbin'
summary(object, prn = TRUE, dec = 3, ...)
\end{verbatim}
\end{Usage}
%
\begin{Arguments}
\begin{ldescription}
\item[\code{object}] Return value from \code{\LinkA{evalbin}{evalbin}}

\item[\code{prn}] Print full table of measures per model and bin

\item[\code{dec}] Number of decimals to show

\item[\code{...}] further arguments passed to or from other methods
\end{ldescription}
\end{Arguments}
%
\begin{Details}
See \url{https://radiant-rstats.github.io/docs/model/evalbin.html} for an example in Radiant
\end{Details}
%
\begin{SeeAlso}
\code{\LinkA{evalbin}{evalbin}} to summarize results

\code{\LinkA{plot.evalbin}{plot.evalbin}} to plot results
\end{SeeAlso}
%
\begin{Examples}
\begin{ExampleCode}
data.frame(buy = dvd$buy, pred1 = runif(20000), pred2 = ifelse(dvd$buy == "yes", 1, 0)) %>%
  evalbin(c("pred1", "pred2"), "buy") %>%
  summary()
\end{ExampleCode}
\end{Examples}
\HeaderA{summary.evalreg}{Summary method for the evalreg function}{summary.evalreg}
%
\begin{Description}
Summary method for the evalreg function
\end{Description}
%
\begin{Usage}
\begin{verbatim}
## S3 method for class 'evalreg'
summary(object, dec = 3, ...)
\end{verbatim}
\end{Usage}
%
\begin{Arguments}
\begin{ldescription}
\item[\code{object}] Return value from \code{\LinkA{evalreg}{evalreg}}

\item[\code{dec}] Number of decimals to show

\item[\code{...}] further arguments passed to or from other methods
\end{ldescription}
\end{Arguments}
%
\begin{Details}
See \url{https://radiant-rstats.github.io/docs/model/evalreg.html} for an example in Radiant
\end{Details}
%
\begin{SeeAlso}
\code{\LinkA{evalreg}{evalreg}} to summarize results

\code{\LinkA{plot.evalreg}{plot.evalreg}} to plot results
\end{SeeAlso}
%
\begin{Examples}
\begin{ExampleCode}
data.frame(price = diamonds$price, pred1 = rnorm(3000), pred2 = diamonds$price) %>%
  evalreg(pred = c("pred1", "pred2"), "price") %>%
  summary()

\end{ExampleCode}
\end{Examples}
\HeaderA{summary.gbt}{Summary method for the gbt function}{summary.gbt}
%
\begin{Description}
Summary method for the gbt function
\end{Description}
%
\begin{Usage}
\begin{verbatim}
## S3 method for class 'gbt'
summary(object, prn = TRUE, ...)
\end{verbatim}
\end{Usage}
%
\begin{Arguments}
\begin{ldescription}
\item[\code{object}] Return value from \code{\LinkA{gbt}{gbt}}

\item[\code{prn}] Print iteration history

\item[\code{...}] further arguments passed to or from other methods
\end{ldescription}
\end{Arguments}
%
\begin{Details}
See \url{https://radiant-rstats.github.io/docs/model/gbt.html} for an example in Radiant
\end{Details}
%
\begin{SeeAlso}
\code{\LinkA{gbt}{gbt}} to generate results

\code{\LinkA{plot.gbt}{plot.gbt}} to plot results

\code{\LinkA{predict.gbt}{predict.gbt}} for prediction
\end{SeeAlso}
%
\begin{Examples}
\begin{ExampleCode}
result <- gbt(
  titanic, "survived", c("pclass", "sex"),
  early_stopping_rounds = 0, nthread = 1
)
summary(result)
\end{ExampleCode}
\end{Examples}
\HeaderA{summary.logistic}{Summary method for the logistic function}{summary.logistic}
%
\begin{Description}
Summary method for the logistic function
\end{Description}
%
\begin{Usage}
\begin{verbatim}
## S3 method for class 'logistic'
summary(object, sum_check = "", conf_lev = 0.95, test_var = "", dec = 3, ...)
\end{verbatim}
\end{Usage}
%
\begin{Arguments}
\begin{ldescription}
\item[\code{object}] Return value from \code{\LinkA{logistic}{logistic}}

\item[\code{sum\_check}] Optional output. "vif" to show multicollinearity diagnostics. "confint" to show coefficient confidence interval estimates. "odds" to show odds ratios and confidence interval estimates.

\item[\code{conf\_lev}] Confidence level to use for coefficient and odds confidence intervals (.95 is the default)

\item[\code{test\_var}] Variables to evaluate in model comparison (i.e., a competing models Chi-squared test)

\item[\code{dec}] Number of decimals to show

\item[\code{...}] further arguments passed to or from other methods
\end{ldescription}
\end{Arguments}
%
\begin{Details}
See \url{https://radiant-rstats.github.io/docs/model/logistic.html} for an example in Radiant
\end{Details}
%
\begin{SeeAlso}
\code{\LinkA{logistic}{logistic}} to generate the results

\code{\LinkA{plot.logistic}{plot.logistic}} to plot the results

\code{\LinkA{predict.logistic}{predict.logistic}} to generate predictions

\code{\LinkA{plot.model.predict}{plot.model.predict}} to plot prediction output
\end{SeeAlso}
%
\begin{Examples}
\begin{ExampleCode}

result <- logistic(titanic, "survived", "pclass", lev = "Yes")
result <- logistic(titanic, "survived", "pclass", lev = "Yes")
summary(result, test_var = "pclass")
res <- logistic(titanic, "survived", c("pclass", "sex"), int = "pclass:sex", lev = "Yes")
summary(res, sum_check = c("vif", "confint", "odds"))
titanic %>%
  logistic("survived", c("pclass", "sex", "age"), lev = "Yes") %>%
  summary("vif")
\end{ExampleCode}
\end{Examples}
\HeaderA{summary.mnl}{Summary method for the mnl function}{summary.mnl}
%
\begin{Description}
Summary method for the mnl function
\end{Description}
%
\begin{Usage}
\begin{verbatim}
## S3 method for class 'mnl'
summary(object, sum_check = "", conf_lev = 0.95, test_var = "", dec = 3, ...)
\end{verbatim}
\end{Usage}
%
\begin{Arguments}
\begin{ldescription}
\item[\code{object}] Return value from \code{\LinkA{mnl}{mnl}}

\item[\code{sum\_check}] Optional output. "confint" to show coefficient confidence interval estimates. "rrr" to show relative risk ratios (RRRs) and confidence interval estimates.

\item[\code{conf\_lev}] Confidence level to use for coefficient and RRRs confidence intervals (.95 is the default)

\item[\code{test\_var}] Variables to evaluate in model comparison (i.e., a competing models Chi-squared test)

\item[\code{dec}] Number of decimals to show

\item[\code{...}] further arguments passed to or from other methods
\end{ldescription}
\end{Arguments}
%
\begin{Details}
See \url{https://radiant-rstats.github.io/docs/model/mnl.html} for an example in Radiant
\end{Details}
%
\begin{SeeAlso}
\code{\LinkA{mnl}{mnl}} to generate the results

\code{\LinkA{plot.mnl}{plot.mnl}} to plot the results

\code{\LinkA{predict.mnl}{predict.mnl}} to generate predictions

\code{\LinkA{plot.model.predict}{plot.model.predict}} to plot prediction output
\end{SeeAlso}
%
\begin{Examples}
\begin{ExampleCode}
result <- mnl(
  ketchup,
  rvar = "choice",
  evar = c("price.heinz28", "price.heinz32", "price.heinz41", "price.hunts32"),
  lev = "heinz28"
)
summary(result)

\end{ExampleCode}
\end{Examples}
\HeaderA{summary.nb}{Summary method for the nb function}{summary.nb}
%
\begin{Description}
Summary method for the nb function
\end{Description}
%
\begin{Usage}
\begin{verbatim}
## S3 method for class 'nb'
summary(object, dec = 3, ...)
\end{verbatim}
\end{Usage}
%
\begin{Arguments}
\begin{ldescription}
\item[\code{object}] Return value from \code{\LinkA{nb}{nb}}

\item[\code{dec}] Decimals

\item[\code{...}] further arguments passed to or from other methods
\end{ldescription}
\end{Arguments}
%
\begin{Details}
See \url{https://radiant-rstats.github.io/docs/model/nb.html} for an example in Radiant
\end{Details}
%
\begin{SeeAlso}
\code{\LinkA{nb}{nb}} to generate results

\code{\LinkA{plot.nb}{plot.nb}} to plot results

\code{\LinkA{predict.nb}{predict.nb}} for prediction
\end{SeeAlso}
%
\begin{Examples}
\begin{ExampleCode}
result <- nb(titanic, "survived", c("pclass", "sex", "age"))
summary(result)

\end{ExampleCode}
\end{Examples}
\HeaderA{summary.nn}{Summary method for the nn function}{summary.nn}
%
\begin{Description}
Summary method for the nn function
\end{Description}
%
\begin{Usage}
\begin{verbatim}
## S3 method for class 'nn'
summary(object, prn = TRUE, ...)
\end{verbatim}
\end{Usage}
%
\begin{Arguments}
\begin{ldescription}
\item[\code{object}] Return value from \code{\LinkA{nn}{nn}}

\item[\code{prn}] Print list of weights

\item[\code{...}] further arguments passed to or from other methods
\end{ldescription}
\end{Arguments}
%
\begin{Details}
See \url{https://radiant-rstats.github.io/docs/model/nn.html} for an example in Radiant
\end{Details}
%
\begin{SeeAlso}
\code{\LinkA{nn}{nn}} to generate results

\code{\LinkA{plot.nn}{plot.nn}} to plot results

\code{\LinkA{predict.nn}{predict.nn}} for prediction
\end{SeeAlso}
%
\begin{Examples}
\begin{ExampleCode}
result <- nn(titanic, "survived", "pclass", lev = "Yes")
summary(result)
\end{ExampleCode}
\end{Examples}
\HeaderA{summary.regress}{Summary method for the regress function}{summary.regress}
%
\begin{Description}
Summary method for the regress function
\end{Description}
%
\begin{Usage}
\begin{verbatim}
## S3 method for class 'regress'
summary(object, sum_check = "", conf_lev = 0.95, test_var = "", dec = 3, ...)
\end{verbatim}
\end{Usage}
%
\begin{Arguments}
\begin{ldescription}
\item[\code{object}] Return value from \code{\LinkA{regress}{regress}}

\item[\code{sum\_check}] Optional output. "rsme" to show the root mean squared error and the standard deviation of the residuals. "sumsquares" to show the sum of squares table. "vif" to show multicollinearity diagnostics. "confint" to show coefficient confidence interval estimates.

\item[\code{conf\_lev}] Confidence level used to estimate confidence intervals (.95 is the default)

\item[\code{test\_var}] Variables to evaluate in model comparison (i.e., a competing models F-test)

\item[\code{dec}] Number of decimals to show

\item[\code{...}] further arguments passed to or from other methods
\end{ldescription}
\end{Arguments}
%
\begin{Details}
See \url{https://radiant-rstats.github.io/docs/model/regress.html} for an example in Radiant
\end{Details}
%
\begin{SeeAlso}
\code{\LinkA{regress}{regress}} to generate the results

\code{\LinkA{plot.regress}{plot.regress}} to plot results

\code{\LinkA{predict.regress}{predict.regress}} to generate predictions
\end{SeeAlso}
%
\begin{Examples}
\begin{ExampleCode}
result <- regress(diamonds, "price", c("carat", "clarity"))
summary(result, sum_check = c("rmse", "sumsquares", "vif", "confint"), test_var = "clarity")
result <- regress(ideal, "y", c("x1", "x2"))
summary(result, test_var = "x2")
ideal %>%
  regress("y", "x1:x3") %>%
  summary()

\end{ExampleCode}
\end{Examples}
\HeaderA{summary.repeater}{Summarize repeated simulation}{summary.repeater}
%
\begin{Description}
Summarize repeated simulation
\end{Description}
%
\begin{Usage}
\begin{verbatim}
## S3 method for class 'repeater'
summary(object, dec = 4, ...)
\end{verbatim}
\end{Usage}
%
\begin{Arguments}
\begin{ldescription}
\item[\code{object}] Return value from \code{\LinkA{repeater}{repeater}}

\item[\code{dec}] Number of decimals to show

\item[\code{...}] further arguments passed to or from other methods
\end{ldescription}
\end{Arguments}
%
\begin{SeeAlso}
\code{\LinkA{repeater}{repeater}} to run a repeated simulation

\code{\LinkA{plot.repeater}{plot.repeater}} to plot results from repeated simulation
\end{SeeAlso}
\HeaderA{summary.rforest}{Summary method for the rforest function}{summary.rforest}
%
\begin{Description}
Summary method for the rforest function
\end{Description}
%
\begin{Usage}
\begin{verbatim}
## S3 method for class 'rforest'
summary(object, ...)
\end{verbatim}
\end{Usage}
%
\begin{Arguments}
\begin{ldescription}
\item[\code{object}] Return value from \code{\LinkA{rforest}{rforest}}

\item[\code{...}] further arguments passed to or from other methods
\end{ldescription}
\end{Arguments}
%
\begin{Details}
See \url{https://radiant-rstats.github.io/docs/model/rforest.html} for an example in Radiant
\end{Details}
%
\begin{SeeAlso}
\code{\LinkA{rforest}{rforest}} to generate results

\code{\LinkA{plot.rforest}{plot.rforest}} to plot results

\code{\LinkA{predict.rforest}{predict.rforest}} for prediction
\end{SeeAlso}
%
\begin{Examples}
\begin{ExampleCode}
result <- rforest(titanic, "survived", "pclass", lev = "Yes")
summary(result)

\end{ExampleCode}
\end{Examples}
\HeaderA{summary.simulater}{Summary method for the simulater function}{summary.simulater}
%
\begin{Description}
Summary method for the simulater function
\end{Description}
%
\begin{Usage}
\begin{verbatim}
## S3 method for class 'simulater'
summary(object, dec = 4, ...)
\end{verbatim}
\end{Usage}
%
\begin{Arguments}
\begin{ldescription}
\item[\code{object}] Return value from \code{\LinkA{simulater}{simulater}}

\item[\code{dec}] Number of decimals to show

\item[\code{...}] further arguments passed to or from other methods
\end{ldescription}
\end{Arguments}
%
\begin{Details}
See \url{https://radiant-rstats.github.io/docs/model/simulater.html} for an example in Radiant
\end{Details}
%
\begin{SeeAlso}
\code{\LinkA{simulater}{simulater}} to generate the results

\code{\LinkA{plot.simulater}{plot.simulater}} to plot results
\end{SeeAlso}
%
\begin{Examples}
\begin{ExampleCode}
simdat <- simulater(norm = "demand 2000 1000", seed = 1234)
summary(simdat)

\end{ExampleCode}
\end{Examples}
\HeaderA{summary.uplift}{Summary method for the uplift function}{summary.uplift}
%
\begin{Description}
Summary method for the uplift function
\end{Description}
%
\begin{Usage}
\begin{verbatim}
## S3 method for class 'uplift'
summary(object, prn = TRUE, dec = 3, ...)
\end{verbatim}
\end{Usage}
%
\begin{Arguments}
\begin{ldescription}
\item[\code{object}] Return value from \code{\LinkA{evalbin}{evalbin}}

\item[\code{prn}] Print full table of measures per model and bin

\item[\code{dec}] Number of decimals to show

\item[\code{...}] further arguments passed to or from other methods
\end{ldescription}
\end{Arguments}
%
\begin{Details}
See \url{https://radiant-rstats.github.io/docs/model/evalbin.html} for an example in Radiant
\end{Details}
%
\begin{SeeAlso}
\code{\LinkA{evalbin}{evalbin}} to summarize results

\code{\LinkA{plot.evalbin}{plot.evalbin}} to plot results
\end{SeeAlso}
%
\begin{Examples}
\begin{ExampleCode}
data.frame(buy = dvd$buy, pred1 = runif(20000), pred2 = ifelse(dvd$buy == "yes", 1, 0)) %>%
  evalbin(c("pred1", "pred2"), "buy") %>%
  summary()
\end{ExampleCode}
\end{Examples}
\HeaderA{test\_specs}{Add interaction terms to list of test variables if needed}{test.Rul.specs}
%
\begin{Description}
Add interaction terms to list of test variables if needed
\end{Description}
%
\begin{Usage}
\begin{verbatim}
test_specs(tv, int)
\end{verbatim}
\end{Usage}
%
\begin{Arguments}
\begin{ldescription}
\item[\code{tv}] List of variables to use for testing for regress or logistic

\item[\code{int}] Interaction terms specified
\end{ldescription}
\end{Arguments}
%
\begin{Details}
See \url{https://radiant-rstats.github.io/docs/model/regress.html} for an example in Radiant
\end{Details}
%
\begin{Value}
A vector of variables names to test
\end{Value}
%
\begin{Examples}
\begin{ExampleCode}
test_specs("a", "a:b")
test_specs("a", c("a:b", "b:c"))
test_specs("a", c("a:b", "b:c", "I(c^2)"))
test_specs(c("a", "b", "c"), c("a:b", "b:c", "I(c^2)"))

\end{ExampleCode}
\end{Examples}
\HeaderA{uplift}{Evaluate uplift for different (binary) classification models}{uplift}
%
\begin{Description}
Evaluate uplift for different (binary) classification models
\end{Description}
%
\begin{Usage}
\begin{verbatim}
uplift(
  dataset,
  pred,
  rvar,
  lev = "",
  tvar,
  tlev = "",
  qnt = 10,
  cost = 1,
  margin = 2,
  scale = 1,
  train = "All",
  data_filter = "",
  arr = "",
  rows = NULL,
  envir = parent.frame()
)
\end{verbatim}
\end{Usage}
%
\begin{Arguments}
\begin{ldescription}
\item[\code{dataset}] Dataset

\item[\code{pred}] Predictions or predictors

\item[\code{rvar}] Response variable

\item[\code{lev}] The level in the response variable defined as success

\item[\code{tvar}] Treatment variable

\item[\code{tlev}] The level in the treatment variable defined as the treatment

\item[\code{qnt}] Number of bins to create

\item[\code{cost}] Cost for each connection (e.g., email or mailing)

\item[\code{margin}] Margin on each customer purchase

\item[\code{scale}] Scaling factor to apply to calculations

\item[\code{train}] Use data from training ("Training"), test ("Test"), both ("Both"), or all data ("All") to evaluate model evalbin

\item[\code{data\_filter}] Expression entered in, e.g., Data > View to filter the dataset in Radiant. The expression should be a string (e.g., "price > 10000")

\item[\code{arr}] Expression to arrange (sort) the data on (e.g., "color, desc(price)")

\item[\code{rows}] Rows to select from the specified dataset

\item[\code{envir}] Environment to extract data from
\end{ldescription}
\end{Arguments}
%
\begin{Details}
Evaluate uplift for different (binary) classification models based on predictions. See \url{https://radiant-rstats.github.io/docs/model/evalbin.html} for an example in Radiant
\end{Details}
%
\begin{Value}
A list of results
\end{Value}
%
\begin{SeeAlso}
\code{\LinkA{summary.evalbin}{summary.evalbin}} to summarize results

\code{\LinkA{plot.evalbin}{plot.evalbin}} to plot results
\end{SeeAlso}
%
\begin{Examples}
\begin{ExampleCode}
data.frame(buy = dvd$buy, pred1 = runif(20000), pred2 = ifelse(dvd$buy == "yes", 1, 0)) %>%
  evalbin(c("pred1", "pred2"), "buy") %>%
  str()
\end{ExampleCode}
\end{Examples}
\HeaderA{varimp}{Variable importance using the vip package and permutation importance}{varimp}
%
\begin{Description}
Variable importance using the vip package and permutation importance
\end{Description}
%
\begin{Usage}
\begin{verbatim}
varimp(object, rvar, lev, data = NULL, seed = 1234)
\end{verbatim}
\end{Usage}
%
\begin{Arguments}
\begin{ldescription}
\item[\code{object}] Model object created by Radiant

\item[\code{rvar}] Label to identify the response or target variable

\item[\code{lev}] Reference class for binary classifier (rvar)

\item[\code{data}] Data to use for prediction. Will default to the data used to estimate the model

\item[\code{seed}] Random seed for reproducibility
\end{ldescription}
\end{Arguments}
\HeaderA{varimp\_plot}{Plot permutation importance}{varimp.Rul.plot}
%
\begin{Description}
Plot permutation importance
\end{Description}
%
\begin{Usage}
\begin{verbatim}
varimp_plot(object, rvar, lev, data = NULL, seed = 1234)
\end{verbatim}
\end{Usage}
%
\begin{Arguments}
\begin{ldescription}
\item[\code{object}] Model object created by Radiant

\item[\code{rvar}] Label to identify the response or target variable

\item[\code{lev}] Reference class for binary classifier (rvar)

\item[\code{data}] Data to use for prediction. Will default to the data used to estimate the model

\item[\code{seed}] Random seed for reproducibility
\end{ldescription}
\end{Arguments}
\HeaderA{var\_check}{Check if main effects for all interaction effects are included in the model}{var.Rul.check}
%
\begin{Description}
Check if main effects for all interaction effects are included in the model
\end{Description}
%
\begin{Usage}
\begin{verbatim}
var_check(ev, cn, intv = c())
\end{verbatim}
\end{Usage}
%
\begin{Arguments}
\begin{ldescription}
\item[\code{ev}] List of explanatory variables provided to \code{\LinkA{regress}{regress}} or \code{\LinkA{logistic}{logistic}}

\item[\code{cn}] Column names for all explanatory variables in the dataset

\item[\code{intv}] Interaction terms specified
\end{ldescription}
\end{Arguments}
%
\begin{Details}
If ':' is used to select a range evar is updated. See \url{https://radiant-rstats.github.io/docs/model/regress.html} for an example in Radiant
\end{Details}
%
\begin{Value}
\code{vars} is a vector of right-hand side variables, possibly with interactions, \code{iv} is the list of explanatory variables, and \code{intv} are interaction terms
\end{Value}
%
\begin{Examples}
\begin{ExampleCode}
var_check("a:d", c("a", "b", "c", "d"))
var_check(c("a", "b"), c("a", "b"), "a:c")
var_check(c("a", "b"), c("a", "b"), "a:c")
var_check(c("a", "b"), c("a", "b"), c("a:c", "I(b^2)"))

\end{ExampleCode}
\end{Examples}
\HeaderA{write.coeff}{Write coefficient table for linear and logistic regression}{write.coeff}
%
\begin{Description}
Write coefficient table for linear and logistic regression
\end{Description}
%
\begin{Usage}
\begin{verbatim}
write.coeff(object, file = "", sort = FALSE, intercept = TRUE)
\end{verbatim}
\end{Usage}
%
\begin{Arguments}
\begin{ldescription}
\item[\code{object}] A fitted model object of class regress or logistic

\item[\code{file}] A character string naming a file. "" indicates output to the console

\item[\code{sort}] Sort table by variable importance

\item[\code{intercept}] Include the intercept in the output (TRUE or FALSE). TRUE is the default
\end{ldescription}
\end{Arguments}
%
\begin{Details}
Write coefficients and importance scores to csv or or return as a data.frame
\end{Details}
%
\begin{Examples}
\begin{ExampleCode}

regress(
  diamonds,
  rvar = "price", evar = c("carat", "clarity", "color", "x"),
  int = c("carat:clarity", "clarity:color", "I(x^2)"), check = "standardize"
) %>%
  write.coeff(sort = TRUE) %>%
  format_df(dec = 3)

logistic(titanic, "survived", c("pclass", "sex"), lev = "Yes") %>%
  write.coeff(intercept = FALSE, sort = TRUE) %>%
  format_df(dec = 2)
\end{ExampleCode}
\end{Examples}
\printindex{}
\end{document}
